% ============================================================================
% CAPÍTULO 3 - FUNDAMENTAÇÃO TEÓRICA
% ============================================================================

\section{Fundamentação Teórica}

% --- Slide: Detectores de Objetos ---
\begin{frame}{Detectores de Objetos}
  \vspace{-2mm}
  {\small\color{sublabelcolor}%
    Evolução dos detectores: \textbf{Two-stage} (alta precisão) vs \textbf{Single-stage} (baixa latência)}

  \vfill
  \begin{center}
    \resizebox{!}{0.65\textheight}{% ============================================================================
% Diagrama Comparativo: Detectores de Dois Estágios vs Um Estágio
% Para uso em dissertação de mestrado - PPGC/UFRGS
% Uso: % ============================================================================
% Diagrama Comparativo: Detectores de Dois Estágios vs Um Estágio
% Para uso em dissertação de mestrado - PPGC/UFRGS
% Uso: % ============================================================================
% Diagrama Comparativo: Detectores de Dois Estágios vs Um Estágio
% Para uso em dissertação de mestrado - PPGC/UFRGS
% Uso: \input{figuras/diagrama_detectores.tex}
% ============================================================================

\begin{tikzpicture}[
    node distance=0.6cm and 0.8cm,
    scale=0.85, transform shape,
    % Estilos dos blocos
    block/.style={
        rectangle,
        rounded corners=3pt,
        minimum width=2.6cm,
        minimum height=0.7cm,
        text centered,
        font=\small\sffamily,
        text=white,
        draw=none,
    },
    inputblock/.style={block, fill=inputcolor, minimum height=0.9cm},
    backboneblock/.style={block, fill=backbonecolor},
    rpnblock/.style={block, fill=rpncolor},
    roiblock/.style={block, fill=roicolor},
    headblock/.style={block, fill=headcolor},
    outputblock/.style={block, fill=outputcolor, minimum height=0.8cm},
    detectionblock/.style={block, fill=detectioncolor},
    gridblock/.style={block, fill=gridcolor},
    % Estilos das setas
    arrow/.style={-{Stealth[length=2mm, width=1.5mm]}, thick, arrowcolor},
    dashedarrow/.style={-{Stealth[length=2mm, width=1.5mm]}, thick, arrowcolor, dashed},
    % Estilos das labels
    titlelabel/.style={font=\normalsize\sffamily\bfseries, text=labelcolor},
    stagelabel/.style={font=\tiny\sffamily\bfseries, text=white, fill=stagecolor, rounded corners=2pt, inner sep=3pt},
    sublabel/.style={font=\tiny\sffamily, text=sublabelcolor},
    divider/.style={thick, dashed, dividercolor},
]

% ============================================================================
% TÍTULO
% ============================================================================
\node[titlelabel] at (0, 0.3) {Comparação de Arquiteturas de Detectores};

% ============================================================================
% DOIS ESTÁGIOS (Faster R-CNN)
% ============================================================================
\begin{scope}[shift={(-4, 0)}]
    \node[titlelabel] at (0, -0.5) {Dois Estágios};
    \node[sublabel] at (0, -0.9) {(Faster R-CNN)};

    \node[stagelabel] at (-1.8, -1.6) {Estágio 1};
    \node[inputblock] (input2) at (0, -1.6) {Imagem};
    \node[backboneblock] (bb2) at (0, -2.7) {Backbone};
    \node[rpnblock] (rpn) at (0, -3.8) {RPN};
    \node[sublabel, below=0.02cm of rpn] {Propostas de regiões};

    \node[stagelabel] at (-1.8, -4.9) {Estágio 2};
    \node[roiblock] (roi) at (0, -5.0) {ROI Pooling};
    \node[headblock] (head2) at (0, -6.1) {Classificação};
    \node[outputblock] (out2) at (0, -7.2) {Detecções};

    \draw[arrow] (input2) -- (bb2);
    \draw[arrow] (bb2) -- (rpn);
    \draw[arrow] (rpn) -- (roi);
    \draw[arrow] (roi) -- (head2);
    \draw[arrow] (head2) -- (out2);
    \draw[dashedarrow] (bb2.east) -- ++(0.5,0) |- (roi.east);

    \node[sublabel] at (0, -8.0) {Alta precisão | Maior latência};
\end{scope}

% ============================================================================
% LINHA DIVISÓRIA
% ============================================================================
\draw[divider] (0, -0.3) -- (0, -8.2);

% ============================================================================
% UM ESTÁGIO (YOLO)
% ============================================================================
\begin{scope}[shift={(4, 0)}]
    \node[titlelabel] at (0, -0.5) {Um Estágio};
    \node[sublabel] at (0, -0.9) {(YOLO, SSD)};

    \node[stagelabel] at (1.8, -1.6) {Único};
    \node[inputblock] (input1) at (0, -1.6) {Imagem};
    \node[backboneblock] (bb1) at (0, -2.7) {Backbone};
    \node[detectionblock] (det) at (0, -3.8) {Detecção Direta};
    \node[sublabel, below=0.02cm of det] {Predição em grid};
    \node[gridblock] (grid) at (0, -5.0) {Grid $S \times S$};
    \node[outputblock] (out1) at (0, -6.1) {Boxes + Classes};

    \draw[arrow] (input1) -- (bb1);
    \draw[arrow] (bb1) -- (det);
    \draw[arrow] (det) -- (grid);
    \draw[arrow] (grid) -- (out1);
    \draw[dashedarrow] (bb1.west) -- ++(-0.5,0) |- (det.west);
    \node[sublabel, rotate=90] at (-1.0, -3.2) {Multi-escala};

    \node[sublabel] at (0, -6.9) {Baixa latência | Eficiente};
\end{scope}

\end{tikzpicture}

% ============================================================================

\begin{tikzpicture}[
    node distance=0.6cm and 0.8cm,
    scale=0.85, transform shape,
    % Estilos dos blocos
    block/.style={
        rectangle,
        rounded corners=3pt,
        minimum width=2.6cm,
        minimum height=0.7cm,
        text centered,
        font=\small\sffamily,
        text=white,
        draw=none,
    },
    inputblock/.style={block, fill=inputcolor, minimum height=0.9cm},
    backboneblock/.style={block, fill=backbonecolor},
    rpnblock/.style={block, fill=rpncolor},
    roiblock/.style={block, fill=roicolor},
    headblock/.style={block, fill=headcolor},
    outputblock/.style={block, fill=outputcolor, minimum height=0.8cm},
    detectionblock/.style={block, fill=detectioncolor},
    gridblock/.style={block, fill=gridcolor},
    % Estilos das setas
    arrow/.style={-{Stealth[length=2mm, width=1.5mm]}, thick, arrowcolor},
    dashedarrow/.style={-{Stealth[length=2mm, width=1.5mm]}, thick, arrowcolor, dashed},
    % Estilos das labels
    titlelabel/.style={font=\normalsize\sffamily\bfseries, text=labelcolor},
    stagelabel/.style={font=\tiny\sffamily\bfseries, text=white, fill=stagecolor, rounded corners=2pt, inner sep=3pt},
    sublabel/.style={font=\tiny\sffamily, text=sublabelcolor},
    divider/.style={thick, dashed, dividercolor},
]

% ============================================================================
% TÍTULO
% ============================================================================
\node[titlelabel] at (0, 0.3) {Comparação de Arquiteturas de Detectores};

% ============================================================================
% DOIS ESTÁGIOS (Faster R-CNN)
% ============================================================================
\begin{scope}[shift={(-4, 0)}]
    \node[titlelabel] at (0, -0.5) {Dois Estágios};
    \node[sublabel] at (0, -0.9) {(Faster R-CNN)};

    \node[stagelabel] at (-1.8, -1.6) {Estágio 1};
    \node[inputblock] (input2) at (0, -1.6) {Imagem};
    \node[backboneblock] (bb2) at (0, -2.7) {Backbone};
    \node[rpnblock] (rpn) at (0, -3.8) {RPN};
    \node[sublabel, below=0.02cm of rpn] {Propostas de regiões};

    \node[stagelabel] at (-1.8, -4.9) {Estágio 2};
    \node[roiblock] (roi) at (0, -5.0) {ROI Pooling};
    \node[headblock] (head2) at (0, -6.1) {Classificação};
    \node[outputblock] (out2) at (0, -7.2) {Detecções};

    \draw[arrow] (input2) -- (bb2);
    \draw[arrow] (bb2) -- (rpn);
    \draw[arrow] (rpn) -- (roi);
    \draw[arrow] (roi) -- (head2);
    \draw[arrow] (head2) -- (out2);
    \draw[dashedarrow] (bb2.east) -- ++(0.5,0) |- (roi.east);

    \node[sublabel] at (0, -8.0) {Alta precisão | Maior latência};
\end{scope}

% ============================================================================
% LINHA DIVISÓRIA
% ============================================================================
\draw[divider] (0, -0.3) -- (0, -8.2);

% ============================================================================
% UM ESTÁGIO (YOLO)
% ============================================================================
\begin{scope}[shift={(4, 0)}]
    \node[titlelabel] at (0, -0.5) {Um Estágio};
    \node[sublabel] at (0, -0.9) {(YOLO, SSD)};

    \node[stagelabel] at (1.8, -1.6) {Único};
    \node[inputblock] (input1) at (0, -1.6) {Imagem};
    \node[backboneblock] (bb1) at (0, -2.7) {Backbone};
    \node[detectionblock] (det) at (0, -3.8) {Detecção Direta};
    \node[sublabel, below=0.02cm of det] {Predição em grid};
    \node[gridblock] (grid) at (0, -5.0) {Grid $S \times S$};
    \node[outputblock] (out1) at (0, -6.1) {Boxes + Classes};

    \draw[arrow] (input1) -- (bb1);
    \draw[arrow] (bb1) -- (det);
    \draw[arrow] (det) -- (grid);
    \draw[arrow] (grid) -- (out1);
    \draw[dashedarrow] (bb1.west) -- ++(-0.5,0) |- (det.west);
    \node[sublabel, rotate=90] at (-1.0, -3.2) {Multi-escala};

    \node[sublabel] at (0, -6.9) {Baixa latência | Eficiente};
\end{scope}

\end{tikzpicture}

% ============================================================================

\begin{tikzpicture}[
    node distance=0.6cm and 0.8cm,
    scale=0.85, transform shape,
    % Estilos dos blocos
    block/.style={
        rectangle,
        rounded corners=3pt,
        minimum width=2.6cm,
        minimum height=0.7cm,
        text centered,
        font=\small\sffamily,
        text=white,
        draw=none,
    },
    inputblock/.style={block, fill=inputcolor, minimum height=0.9cm},
    backboneblock/.style={block, fill=backbonecolor},
    rpnblock/.style={block, fill=rpncolor},
    roiblock/.style={block, fill=roicolor},
    headblock/.style={block, fill=headcolor},
    outputblock/.style={block, fill=outputcolor, minimum height=0.8cm},
    detectionblock/.style={block, fill=detectioncolor},
    gridblock/.style={block, fill=gridcolor},
    % Estilos das setas
    arrow/.style={-{Stealth[length=2mm, width=1.5mm]}, thick, arrowcolor},
    dashedarrow/.style={-{Stealth[length=2mm, width=1.5mm]}, thick, arrowcolor, dashed},
    % Estilos das labels
    titlelabel/.style={font=\normalsize\sffamily\bfseries, text=labelcolor},
    stagelabel/.style={font=\tiny\sffamily\bfseries, text=white, fill=stagecolor, rounded corners=2pt, inner sep=3pt},
    sublabel/.style={font=\tiny\sffamily, text=sublabelcolor},
    divider/.style={thick, dashed, dividercolor},
]

% ============================================================================
% TÍTULO
% ============================================================================
\node[titlelabel] at (0, 0.3) {Comparação de Arquiteturas de Detectores};

% ============================================================================
% DOIS ESTÁGIOS (Faster R-CNN)
% ============================================================================
\begin{scope}[shift={(-4, 0)}]
    \node[titlelabel] at (0, -0.5) {Dois Estágios};
    \node[sublabel] at (0, -0.9) {(Faster R-CNN)};

    \node[stagelabel] at (-1.8, -1.6) {Estágio 1};
    \node[inputblock] (input2) at (0, -1.6) {Imagem};
    \node[backboneblock] (bb2) at (0, -2.7) {Backbone};
    \node[rpnblock] (rpn) at (0, -3.8) {RPN};
    \node[sublabel, below=0.02cm of rpn] {Propostas de regiões};

    \node[stagelabel] at (-1.8, -4.9) {Estágio 2};
    \node[roiblock] (roi) at (0, -5.0) {ROI Pooling};
    \node[headblock] (head2) at (0, -6.1) {Classificação};
    \node[outputblock] (out2) at (0, -7.2) {Detecções};

    \draw[arrow] (input2) -- (bb2);
    \draw[arrow] (bb2) -- (rpn);
    \draw[arrow] (rpn) -- (roi);
    \draw[arrow] (roi) -- (head2);
    \draw[arrow] (head2) -- (out2);
    \draw[dashedarrow] (bb2.east) -- ++(0.5,0) |- (roi.east);

    \node[sublabel] at (0, -8.0) {Alta precisão | Maior latência};
\end{scope}

% ============================================================================
% LINHA DIVISÓRIA
% ============================================================================
\draw[divider] (0, -0.3) -- (0, -8.2);

% ============================================================================
% UM ESTÁGIO (YOLO)
% ============================================================================
\begin{scope}[shift={(4, 0)}]
    \node[titlelabel] at (0, -0.5) {Um Estágio};
    \node[sublabel] at (0, -0.9) {(YOLO, SSD)};

    \node[stagelabel] at (1.8, -1.6) {Único};
    \node[inputblock] (input1) at (0, -1.6) {Imagem};
    \node[backboneblock] (bb1) at (0, -2.7) {Backbone};
    \node[detectionblock] (det) at (0, -3.8) {Detecção Direta};
    \node[sublabel, below=0.02cm of det] {Predição em grid};
    \node[gridblock] (grid) at (0, -5.0) {Grid $S \times S$};
    \node[outputblock] (out1) at (0, -6.1) {Boxes + Classes};

    \draw[arrow] (input1) -- (bb1);
    \draw[arrow] (bb1) -- (det);
    \draw[arrow] (det) -- (grid);
    \draw[arrow] (grid) -- (out1);
    \draw[dashedarrow] (bb1.west) -- ++(-0.5,0) |- (det.west);
    \node[sublabel, rotate=90] at (-1.0, -3.2) {Multi-escala};

    \node[sublabel] at (0, -6.9) {Baixa latência | Eficiente};
\end{scope}

\end{tikzpicture}
}
  \end{center}
  \vfill

  \begin{beamercolorbox}[wd=\textwidth, rounded=true, shadow=false]{block body}
    \centering
    {\footnotesize\color{stagecolor}%
      \textbf{YOLO} escolhido: melhor \textit{trade-off} velocidade/acurácia para inferência em tempo real na borda}
  \end{beamercolorbox}
\end{frame}

% --- Slide: Arquitetura YOLOv8 ---
\begin{frame}{Arquitetura YOLOv8}
  \vspace{-2mm}
  {\small\color{sublabelcolor}%
    Três componentes: \textbf{Backbone} (extração) + \textbf{Neck} (fusão multi-escala) + \textbf{Decoupled Head} (detecção)}

  \vfill
  \begin{center}
    \resizebox{0.92\textwidth}{!}{% ============================================================================
% Diagrama da Arquitetura YOLOv8s (Simplificado)
% Para uso em dissertação de mestrado - PPGC/UFRGS
% Uso: % ============================================================================
% Diagrama da Arquitetura YOLOv8s (Simplificado)
% Para uso em dissertação de mestrado - PPGC/UFRGS
% Uso: % ============================================================================
% Diagrama da Arquitetura YOLOv8s (Simplificado)
% Para uso em dissertação de mestrado - PPGC/UFRGS
% Uso: \input{figuras/diagrama_yolov8.tex}
% ============================================================================

\begin{tikzpicture}[
    node distance=0.5cm and 0.6cm,
    scale=0.9, transform shape,
    % Estilos dos blocos
    block/.style={
        rectangle,
        rounded corners=3pt,
        minimum width=1.6cm,
        minimum height=0.6cm,
        align=center,
        font=\small\sffamily,
        text=white,
        draw=none,
    },
    inputblock/.style={block, fill=inputcolor, minimum width=1.8cm, minimum height=0.8cm},
    cbsblock/.style={block, fill=backbonecolor, minimum width=1.2cm},
    c2fblock/.style={block, fill=c3color, minimum width=1.4cm},
    sppfblock/.style={block, fill=backbonecolor!80!black, minimum width=1.4cm},
    neckblock/.style={block, fill=neckcolor, minimum width=1.4cm},
    headblock/.style={block, fill=headcolor, minimum width=1.8cm, minimum height=0.7cm},
    outputblock/.style={block, fill=outputcolor, minimum width=2.2cm, minimum height=0.8cm},
    % Estilos das setas
    arrow/.style={-{Stealth[length=1.8mm, width=1.2mm]}, thick, arrowcolor},
    skiparrow/.style={-{Stealth[length=1.8mm, width=1.2mm]}, thick, skipcolor},
    % Estilos das labels
    sectionlabel/.style={font=\small\sffamily\bfseries, text=labelcolor},
    sizelabel/.style={font=\scriptsize\sffamily, text=sublabelcolor},
]

% ============================================================================
% ENTRADA
% ============================================================================
\node[inputblock] (input) {Imagem\\640$\times$640};
\node[sectionlabel, above=0.3cm of input] {Entrada};

% ============================================================================
% BACKBONE
% ============================================================================
\node[cbsblock, right=0.8cm of input] (cbs1) {CBS};
\node[c2fblock, right=0.4cm of cbs1] (c2f1) {C2f};
\node[sizelabel, below=0.02cm of c2f1] {\tiny P3};

\node[cbsblock, right=0.4cm of c2f1] (cbs2) {CBS};
\node[c2fblock, right=0.4cm of cbs2] (c2f2) {C2f};
\node[sizelabel, below=0.02cm of c2f2] {\tiny P4};

\node[cbsblock, right=0.4cm of c2f2] (cbs3) {CBS};
\node[c2fblock, right=0.4cm of cbs3] (c2f3) {C2f};
\node[sppfblock, right=0.4cm of c2f3] (sppf) {SPPF};
\node[sizelabel, below=0.02cm of sppf] {\tiny P5};

% Backbone label
\node[sectionlabel] at ($(cbs2)!0.5!(c2f3) + (0, 0.8)$) {Backbone (CSPDarknet)};

% ============================================================================
% NECK (FPN + PANet)
% ============================================================================
\node[neckblock, below=1.0cm of c2f2] (neck1) {FPN};
\node[neckblock, right=0.6cm of neck1] (neck2) {PANet};
\node[sectionlabel] at ($(neck1)!0.5!(neck2) + (0, 0.7)$) {Neck};

% ============================================================================
% HEAD
% ============================================================================
\node[headblock, below=1.0cm of neck1] (head1) {Head\\80$\times$80};
\node[headblock, right=0.5cm of head1] (head2) {Head\\40$\times$40};
\node[headblock, right=0.5cm of head2] (head3) {Head\\20$\times$20};
\node[sectionlabel] at ($(head2) + (0, 0.7)$) {Decoupled Head};

\node[sizelabel, below=0.1cm of head1] {Pequenos};
\node[sizelabel, below=0.1cm of head2] {Médios};
\node[sizelabel, below=0.1cm of head3] {Grandes};

% ============================================================================
% SAÍDA
% ============================================================================
\node[outputblock, below=1.0cm of head2] (output) {Bounding Boxes\\Classes + Confiança};
\node[sectionlabel, below=0.3cm of output] {Saída};

% ============================================================================
% SETAS
% ============================================================================
% Input -> Backbone
\draw[arrow] (input) -- (cbs1);
\draw[arrow] (cbs1) -- (c2f1);
\draw[arrow] (c2f1) -- (cbs2);
\draw[arrow] (cbs2) -- (c2f2);
\draw[arrow] (c2f2) -- (cbs3);
\draw[arrow] (cbs3) -- (c2f3);
\draw[arrow] (c2f3) -- (sppf);

% Backbone -> Neck (skip connections)
\draw[skiparrow] (c2f1.south) -- ++(0,-0.3) -| (neck1.north west);
\draw[skiparrow] (c2f2.south) -- (neck1.north);
\draw[skiparrow] (sppf.south) -- ++(0,-0.2) -| (neck2.north);

% Neck -> Head
\draw[arrow] (neck1.south) -- ++(0,-0.2) -| (head1.north);
\draw[arrow] (neck1.south east) -- ++(0,-0.2) -| (head2.north);
\draw[arrow] (neck2.south) -- ++(0,-0.2) -| (head3.north);

% Head -> Output
\draw[arrow] (head1.south) -- ++(0,-0.2) -| ($(output.north) + (-0.5,0)$);
\draw[arrow] (head2.south) -- (output.north);
\draw[arrow] (head3.south) -- ++(0,-0.2) -| ($(output.north) + (0.5,0)$);

% ============================================================================
% LEGENDA
% ============================================================================
\begin{scope}[shift={($(output) + (4.5, 2.5)$)}]
    \node[font=\small\sffamily\bfseries, text=labelcolor] at (0, 1.2) {Legenda:};

    \node[block, fill=backbonecolor, minimum width=1.0cm, minimum height=0.4cm] at (0, 0.8) {\scriptsize CBS};
    \node[font=\scriptsize\sffamily, text=labelcolor, anchor=west] at (0.6, 0.8) {Conv+BN+SiLU};

    \node[block, fill=c3color, minimum width=1.0cm, minimum height=0.4cm] at (0, 0.4) {\scriptsize C2f};
    \node[font=\scriptsize\sffamily, text=labelcolor, anchor=west] at (0.6, 0.4) {CSP Bottleneck};

    \node[block, fill=headcolor, minimum width=1.0cm, minimum height=0.4cm] at (0, 0.0) {\scriptsize Head};
    \node[font=\scriptsize\sffamily, text=labelcolor, anchor=west] at (0.6, 0.0) {Cls + Reg separados};

    \draw[arrow] (-0.4, -0.4) -- (0.2, -0.4);
    \node[font=\scriptsize\sffamily, text=labelcolor, anchor=west] at (0.4, -0.4) {Fluxo principal};

    \draw[skiparrow] (-0.4, -0.7) -- (0.2, -0.7);
    \node[font=\scriptsize\sffamily, text=labelcolor, anchor=west] at (0.4, -0.7) {Skip connection};
\end{scope}

\end{tikzpicture}

% ============================================================================

\begin{tikzpicture}[
    node distance=0.5cm and 0.6cm,
    scale=0.9, transform shape,
    % Estilos dos blocos
    block/.style={
        rectangle,
        rounded corners=3pt,
        minimum width=1.6cm,
        minimum height=0.6cm,
        align=center,
        font=\small\sffamily,
        text=white,
        draw=none,
    },
    inputblock/.style={block, fill=inputcolor, minimum width=1.8cm, minimum height=0.8cm},
    cbsblock/.style={block, fill=backbonecolor, minimum width=1.2cm},
    c2fblock/.style={block, fill=c3color, minimum width=1.4cm},
    sppfblock/.style={block, fill=backbonecolor!80!black, minimum width=1.4cm},
    neckblock/.style={block, fill=neckcolor, minimum width=1.4cm},
    headblock/.style={block, fill=headcolor, minimum width=1.8cm, minimum height=0.7cm},
    outputblock/.style={block, fill=outputcolor, minimum width=2.2cm, minimum height=0.8cm},
    % Estilos das setas
    arrow/.style={-{Stealth[length=1.8mm, width=1.2mm]}, thick, arrowcolor},
    skiparrow/.style={-{Stealth[length=1.8mm, width=1.2mm]}, thick, skipcolor},
    % Estilos das labels
    sectionlabel/.style={font=\small\sffamily\bfseries, text=labelcolor},
    sizelabel/.style={font=\scriptsize\sffamily, text=sublabelcolor},
]

% ============================================================================
% ENTRADA
% ============================================================================
\node[inputblock] (input) {Imagem\\640$\times$640};
\node[sectionlabel, above=0.3cm of input] {Entrada};

% ============================================================================
% BACKBONE
% ============================================================================
\node[cbsblock, right=0.8cm of input] (cbs1) {CBS};
\node[c2fblock, right=0.4cm of cbs1] (c2f1) {C2f};
\node[sizelabel, below=0.02cm of c2f1] {\tiny P3};

\node[cbsblock, right=0.4cm of c2f1] (cbs2) {CBS};
\node[c2fblock, right=0.4cm of cbs2] (c2f2) {C2f};
\node[sizelabel, below=0.02cm of c2f2] {\tiny P4};

\node[cbsblock, right=0.4cm of c2f2] (cbs3) {CBS};
\node[c2fblock, right=0.4cm of cbs3] (c2f3) {C2f};
\node[sppfblock, right=0.4cm of c2f3] (sppf) {SPPF};
\node[sizelabel, below=0.02cm of sppf] {\tiny P5};

% Backbone label
\node[sectionlabel] at ($(cbs2)!0.5!(c2f3) + (0, 0.8)$) {Backbone (CSPDarknet)};

% ============================================================================
% NECK (FPN + PANet)
% ============================================================================
\node[neckblock, below=1.0cm of c2f2] (neck1) {FPN};
\node[neckblock, right=0.6cm of neck1] (neck2) {PANet};
\node[sectionlabel] at ($(neck1)!0.5!(neck2) + (0, 0.7)$) {Neck};

% ============================================================================
% HEAD
% ============================================================================
\node[headblock, below=1.0cm of neck1] (head1) {Head\\80$\times$80};
\node[headblock, right=0.5cm of head1] (head2) {Head\\40$\times$40};
\node[headblock, right=0.5cm of head2] (head3) {Head\\20$\times$20};
\node[sectionlabel] at ($(head2) + (0, 0.7)$) {Decoupled Head};

\node[sizelabel, below=0.1cm of head1] {Pequenos};
\node[sizelabel, below=0.1cm of head2] {Médios};
\node[sizelabel, below=0.1cm of head3] {Grandes};

% ============================================================================
% SAÍDA
% ============================================================================
\node[outputblock, below=1.0cm of head2] (output) {Bounding Boxes\\Classes + Confiança};
\node[sectionlabel, below=0.3cm of output] {Saída};

% ============================================================================
% SETAS
% ============================================================================
% Input -> Backbone
\draw[arrow] (input) -- (cbs1);
\draw[arrow] (cbs1) -- (c2f1);
\draw[arrow] (c2f1) -- (cbs2);
\draw[arrow] (cbs2) -- (c2f2);
\draw[arrow] (c2f2) -- (cbs3);
\draw[arrow] (cbs3) -- (c2f3);
\draw[arrow] (c2f3) -- (sppf);

% Backbone -> Neck (skip connections)
\draw[skiparrow] (c2f1.south) -- ++(0,-0.3) -| (neck1.north west);
\draw[skiparrow] (c2f2.south) -- (neck1.north);
\draw[skiparrow] (sppf.south) -- ++(0,-0.2) -| (neck2.north);

% Neck -> Head
\draw[arrow] (neck1.south) -- ++(0,-0.2) -| (head1.north);
\draw[arrow] (neck1.south east) -- ++(0,-0.2) -| (head2.north);
\draw[arrow] (neck2.south) -- ++(0,-0.2) -| (head3.north);

% Head -> Output
\draw[arrow] (head1.south) -- ++(0,-0.2) -| ($(output.north) + (-0.5,0)$);
\draw[arrow] (head2.south) -- (output.north);
\draw[arrow] (head3.south) -- ++(0,-0.2) -| ($(output.north) + (0.5,0)$);

% ============================================================================
% LEGENDA
% ============================================================================
\begin{scope}[shift={($(output) + (4.5, 2.5)$)}]
    \node[font=\small\sffamily\bfseries, text=labelcolor] at (0, 1.2) {Legenda:};

    \node[block, fill=backbonecolor, minimum width=1.0cm, minimum height=0.4cm] at (0, 0.8) {\scriptsize CBS};
    \node[font=\scriptsize\sffamily, text=labelcolor, anchor=west] at (0.6, 0.8) {Conv+BN+SiLU};

    \node[block, fill=c3color, minimum width=1.0cm, minimum height=0.4cm] at (0, 0.4) {\scriptsize C2f};
    \node[font=\scriptsize\sffamily, text=labelcolor, anchor=west] at (0.6, 0.4) {CSP Bottleneck};

    \node[block, fill=headcolor, minimum width=1.0cm, minimum height=0.4cm] at (0, 0.0) {\scriptsize Head};
    \node[font=\scriptsize\sffamily, text=labelcolor, anchor=west] at (0.6, 0.0) {Cls + Reg separados};

    \draw[arrow] (-0.4, -0.4) -- (0.2, -0.4);
    \node[font=\scriptsize\sffamily, text=labelcolor, anchor=west] at (0.4, -0.4) {Fluxo principal};

    \draw[skiparrow] (-0.4, -0.7) -- (0.2, -0.7);
    \node[font=\scriptsize\sffamily, text=labelcolor, anchor=west] at (0.4, -0.7) {Skip connection};
\end{scope}

\end{tikzpicture}

% ============================================================================

\begin{tikzpicture}[
    node distance=0.5cm and 0.6cm,
    scale=0.9, transform shape,
    % Estilos dos blocos
    block/.style={
        rectangle,
        rounded corners=3pt,
        minimum width=1.6cm,
        minimum height=0.6cm,
        align=center,
        font=\small\sffamily,
        text=white,
        draw=none,
    },
    inputblock/.style={block, fill=inputcolor, minimum width=1.8cm, minimum height=0.8cm},
    cbsblock/.style={block, fill=backbonecolor, minimum width=1.2cm},
    c2fblock/.style={block, fill=c3color, minimum width=1.4cm},
    sppfblock/.style={block, fill=backbonecolor!80!black, minimum width=1.4cm},
    neckblock/.style={block, fill=neckcolor, minimum width=1.4cm},
    headblock/.style={block, fill=headcolor, minimum width=1.8cm, minimum height=0.7cm},
    outputblock/.style={block, fill=outputcolor, minimum width=2.2cm, minimum height=0.8cm},
    % Estilos das setas
    arrow/.style={-{Stealth[length=1.8mm, width=1.2mm]}, thick, arrowcolor},
    skiparrow/.style={-{Stealth[length=1.8mm, width=1.2mm]}, thick, skipcolor},
    % Estilos das labels
    sectionlabel/.style={font=\small\sffamily\bfseries, text=labelcolor},
    sizelabel/.style={font=\scriptsize\sffamily, text=sublabelcolor},
]

% ============================================================================
% ENTRADA
% ============================================================================
\node[inputblock] (input) {Imagem\\640$\times$640};
\node[sectionlabel, above=0.3cm of input] {Entrada};

% ============================================================================
% BACKBONE
% ============================================================================
\node[cbsblock, right=0.8cm of input] (cbs1) {CBS};
\node[c2fblock, right=0.4cm of cbs1] (c2f1) {C2f};
\node[sizelabel, below=0.02cm of c2f1] {\tiny P3};

\node[cbsblock, right=0.4cm of c2f1] (cbs2) {CBS};
\node[c2fblock, right=0.4cm of cbs2] (c2f2) {C2f};
\node[sizelabel, below=0.02cm of c2f2] {\tiny P4};

\node[cbsblock, right=0.4cm of c2f2] (cbs3) {CBS};
\node[c2fblock, right=0.4cm of cbs3] (c2f3) {C2f};
\node[sppfblock, right=0.4cm of c2f3] (sppf) {SPPF};
\node[sizelabel, below=0.02cm of sppf] {\tiny P5};

% Backbone label
\node[sectionlabel] at ($(cbs2)!0.5!(c2f3) + (0, 0.8)$) {Backbone (CSPDarknet)};

% ============================================================================
% NECK (FPN + PANet)
% ============================================================================
\node[neckblock, below=1.0cm of c2f2] (neck1) {FPN};
\node[neckblock, right=0.6cm of neck1] (neck2) {PANet};
\node[sectionlabel] at ($(neck1)!0.5!(neck2) + (0, 0.7)$) {Neck};

% ============================================================================
% HEAD
% ============================================================================
\node[headblock, below=1.0cm of neck1] (head1) {Head\\80$\times$80};
\node[headblock, right=0.5cm of head1] (head2) {Head\\40$\times$40};
\node[headblock, right=0.5cm of head2] (head3) {Head\\20$\times$20};
\node[sectionlabel] at ($(head2) + (0, 0.7)$) {Decoupled Head};

\node[sizelabel, below=0.1cm of head1] {Pequenos};
\node[sizelabel, below=0.1cm of head2] {Médios};
\node[sizelabel, below=0.1cm of head3] {Grandes};

% ============================================================================
% SAÍDA
% ============================================================================
\node[outputblock, below=1.0cm of head2] (output) {Bounding Boxes\\Classes + Confiança};
\node[sectionlabel, below=0.3cm of output] {Saída};

% ============================================================================
% SETAS
% ============================================================================
% Input -> Backbone
\draw[arrow] (input) -- (cbs1);
\draw[arrow] (cbs1) -- (c2f1);
\draw[arrow] (c2f1) -- (cbs2);
\draw[arrow] (cbs2) -- (c2f2);
\draw[arrow] (c2f2) -- (cbs3);
\draw[arrow] (cbs3) -- (c2f3);
\draw[arrow] (c2f3) -- (sppf);

% Backbone -> Neck (skip connections)
\draw[skiparrow] (c2f1.south) -- ++(0,-0.3) -| (neck1.north west);
\draw[skiparrow] (c2f2.south) -- (neck1.north);
\draw[skiparrow] (sppf.south) -- ++(0,-0.2) -| (neck2.north);

% Neck -> Head
\draw[arrow] (neck1.south) -- ++(0,-0.2) -| (head1.north);
\draw[arrow] (neck1.south east) -- ++(0,-0.2) -| (head2.north);
\draw[arrow] (neck2.south) -- ++(0,-0.2) -| (head3.north);

% Head -> Output
\draw[arrow] (head1.south) -- ++(0,-0.2) -| ($(output.north) + (-0.5,0)$);
\draw[arrow] (head2.south) -- (output.north);
\draw[arrow] (head3.south) -- ++(0,-0.2) -| ($(output.north) + (0.5,0)$);

% ============================================================================
% LEGENDA
% ============================================================================
\begin{scope}[shift={($(output) + (4.5, 2.5)$)}]
    \node[font=\small\sffamily\bfseries, text=labelcolor] at (0, 1.2) {Legenda:};

    \node[block, fill=backbonecolor, minimum width=1.0cm, minimum height=0.4cm] at (0, 0.8) {\scriptsize CBS};
    \node[font=\scriptsize\sffamily, text=labelcolor, anchor=west] at (0.6, 0.8) {Conv+BN+SiLU};

    \node[block, fill=c3color, minimum width=1.0cm, minimum height=0.4cm] at (0, 0.4) {\scriptsize C2f};
    \node[font=\scriptsize\sffamily, text=labelcolor, anchor=west] at (0.6, 0.4) {CSP Bottleneck};

    \node[block, fill=headcolor, minimum width=1.0cm, minimum height=0.4cm] at (0, 0.0) {\scriptsize Head};
    \node[font=\scriptsize\sffamily, text=labelcolor, anchor=west] at (0.6, 0.0) {Cls + Reg separados};

    \draw[arrow] (-0.4, -0.4) -- (0.2, -0.4);
    \node[font=\scriptsize\sffamily, text=labelcolor, anchor=west] at (0.4, -0.4) {Fluxo principal};

    \draw[skiparrow] (-0.4, -0.7) -- (0.2, -0.7);
    \node[font=\scriptsize\sffamily, text=labelcolor, anchor=west] at (0.4, -0.7) {Skip connection};
\end{scope}

\end{tikzpicture}
}
  \end{center}
  \vfill

  \begin{beamercolorbox}[wd=\textwidth, rounded=true, shadow=false]{block body}
    \centering
    {\footnotesize\color{stagecolor}%
      Variantes \textbf{YOLOv8s} e \textbf{YOLOv8m} utilizadas --- \textit{anchor-free} com cabeças desacopladas (cls + reg)}
  \end{beamercolorbox}
\end{frame}

% --- Slide: Transfer Learning ---
\begin{frame}{Transfer Learning}
  \vspace{-2mm}
  {\small\color{sublabelcolor}%
    Transferência de conhecimento: \textbf{COCO} (milhões de imagens) $\rightarrow$ \textbf{domínio específico} (EPIs)}

  \vfill
  \begin{center}
    \resizebox{0.95\textwidth}{!}{% ============================================================================
% Diagrama de Transfer Learning (Simplificado)
% Para uso em dissertacao de mestrado - PPGC/UFRGS
% Uso: % ============================================================================
% Diagrama de Transfer Learning (Simplificado)
% Para uso em dissertacao de mestrado - PPGC/UFRGS
% Uso: % ============================================================================
% Diagrama de Transfer Learning (Simplificado)
% Para uso em dissertacao de mestrado - PPGC/UFRGS
% Uso: \input{figuras/diagrama_transfer_learning.tex}
% ============================================================================

\begin{tikzpicture}[
    scale=0.85, transform shape,
    % Estilos dos blocos
    layerfrozen/.style={
        rectangle, rounded corners=2pt,
        minimum width=2.0cm, minimum height=0.45cm,
        text centered, font=\tiny\sffamily,
        fill=cspcolor, text=white,
        draw=cspcolor!80!black, line width=0.5pt
    },
    layerfinetune/.style={
        rectangle, rounded corners=2pt,
        minimum width=2.0cm, minimum height=0.45cm,
        text centered, font=\tiny\sffamily,
        fill=headcolor, text=white,
        draw=headcolor!80!black, line width=0.5pt
    },
    layernew/.style={
        rectangle, rounded corners=2pt,
        minimum width=2.0cm, minimum height=0.45cm,
        text centered, font=\tiny\sffamily,
        fill=outputcolor, text=white,
        draw=outputcolor!80!black, line width=0.5pt
    },
    datasetbox/.style={
        rectangle, rounded corners=4pt,
        minimum width=2.5cm, minimum height=1.4cm,
        text centered, font=\scriptsize\sffamily,
        fill=#1!10, draw=#1, line width=1pt
    },
    arrow/.style={-{Stealth[length=1.8mm, width=1.2mm]}, thick, arrowcolor},
    transferarrow/.style={-{Stealth[length=2.5mm, width=1.8mm]}, very thick, skipcolor},
    sectionlabel/.style={font=\small\sffamily\bfseries, text=labelcolor},
    sublabel/.style={font=\tiny\sffamily, text=labelcolor}
]

% ============================================================================
% DOMINIO FONTE (Esquerda)
% ============================================================================
\node[sectionlabel] at (0, 3.2) {Domínio Fonte};
\node[sublabel] at (0, 2.8) {(Pré-treinamento)};

\node[datasetbox=inputcolor] (source_data) at (0, 1.5) {};
\node[font=\scriptsize\sffamily\bfseries, text=inputcolor!80!black] at (0, 1.8) {COCO};
\node[font=\tiny\sffamily, text=labelcolor] at (0, 1.2) {Milhões de imagens};

% Camadas fonte (compactas)
\node[layerfrozen] (src_conv1) at (0, -0.1) {Conv - Bordas};
\node[layerfrozen] (src_conv2) at (0, -0.7) {Conv - Texturas};
\node[layerfrozen] (src_conv3) at (0, -1.3) {Conv - Padrões};
\node[layerfinetune] (src_fc1) at (0, -1.9) {FC - Semântica};
\node[layernew] (src_fc2) at (0, -2.5) {FC - Classes};

\draw[arrow] (source_data.south) -- (src_conv1.north);
\draw[arrow] (src_conv1) -- (src_conv2);
\draw[arrow] (src_conv2) -- (src_conv3);
\draw[arrow] (src_conv3) -- (src_fc1);
\draw[arrow] (src_fc1) -- (src_fc2);

\node[sublabel] at (0, -3.1) {80 classes genéricas};

% ============================================================================
% PROCESSO DE TRANSFERENCIA (Centro)
% ============================================================================
\draw[transferarrow] (1.8, -1.0) -- node[above, font=\scriptsize\sffamily\bfseries, text=skipcolor] {Transfer} node[below, font=\tiny\sffamily, text=skipcolor] {Learning} (4.2, -1.0);

% ============================================================================
% DOMINIO ALVO (Direita)
% ============================================================================
\node[sectionlabel] at (6.0, 3.2) {Domínio Alvo};
\node[sublabel] at (6.0, 2.8) {(Fine-tuning)};

\node[datasetbox=headcolor] (target_data) at (6.0, 1.5) {};
\node[font=\scriptsize\sffamily\bfseries, text=headcolor!80!black] at (6.0, 1.8) {Dataset EPI};
\node[font=\tiny\sffamily, text=labelcolor] at (6.0, 1.2) {Classes específicas};

% Camadas alvo
\node[layerfrozen] (tgt_conv1) at (6.0, -0.1) {Conv - Bordas};
\node[layerfrozen] (tgt_conv2) at (6.0, -0.7) {Conv - Texturas};
\node[layerfinetune] (tgt_conv3) at (6.0, -1.3) {Conv - EPIs};
\node[layerfinetune] (tgt_fc1) at (6.0, -1.9) {FC - Atributos};
\node[layernew] (tgt_fc2) at (6.0, -2.5) {FC - Classes EPI};

\draw[arrow] (target_data.south) -- (tgt_conv1.north);
\draw[arrow] (tgt_conv1) -- (tgt_conv2);
\draw[arrow] (tgt_conv2) -- (tgt_conv3);
\draw[arrow] (tgt_conv3) -- (tgt_fc1);
\draw[arrow] (tgt_fc1) -- (tgt_fc2);

\node[sublabel] at (6.0, -3.1) {Capacete, Colete, Luvas...};

% ============================================================================
% LEGENDA COMPACTA
% ============================================================================
\begin{scope}[shift={(9.0, 0.5)}]
    \node[font=\footnotesize\sffamily\bfseries, text=labelcolor] at (0, 0.8) {Legenda:};

    \node[layerfrozen, minimum width=1.2cm] at (0, 0.3) {};
    \node[font=\tiny\sffamily, text=labelcolor, anchor=west] at (0.7, 0.3) {Congelada};

    \node[layerfinetune, minimum width=1.2cm] at (0, -0.2) {};
    \node[font=\tiny\sffamily, text=labelcolor, anchor=west] at (0.7, -0.2) {Ajustada};

    \node[layernew, minimum width=1.2cm] at (0, -0.7) {};
    \node[font=\tiny\sffamily, text=labelcolor, anchor=west] at (0.7, -0.7) {Nova};
\end{scope}

% Seta de conhecimento transferido
\draw[dashed, skipcolor, line width=0.6pt, -{Stealth[length=1.5mm]}]
    (0.8, -0.4) .. controls (2.5, 0.8) and (3.5, 0.8) .. (5.2, -0.4);
\node[font=\tiny\sffamily\itshape, text=skipcolor] at (3.0, 1.2) {Pesos transferidos};

\end{tikzpicture}

% ============================================================================

\begin{tikzpicture}[
    scale=0.85, transform shape,
    % Estilos dos blocos
    layerfrozen/.style={
        rectangle, rounded corners=2pt,
        minimum width=2.0cm, minimum height=0.45cm,
        text centered, font=\tiny\sffamily,
        fill=cspcolor, text=white,
        draw=cspcolor!80!black, line width=0.5pt
    },
    layerfinetune/.style={
        rectangle, rounded corners=2pt,
        minimum width=2.0cm, minimum height=0.45cm,
        text centered, font=\tiny\sffamily,
        fill=headcolor, text=white,
        draw=headcolor!80!black, line width=0.5pt
    },
    layernew/.style={
        rectangle, rounded corners=2pt,
        minimum width=2.0cm, minimum height=0.45cm,
        text centered, font=\tiny\sffamily,
        fill=outputcolor, text=white,
        draw=outputcolor!80!black, line width=0.5pt
    },
    datasetbox/.style={
        rectangle, rounded corners=4pt,
        minimum width=2.5cm, minimum height=1.4cm,
        text centered, font=\scriptsize\sffamily,
        fill=#1!10, draw=#1, line width=1pt
    },
    arrow/.style={-{Stealth[length=1.8mm, width=1.2mm]}, thick, arrowcolor},
    transferarrow/.style={-{Stealth[length=2.5mm, width=1.8mm]}, very thick, skipcolor},
    sectionlabel/.style={font=\small\sffamily\bfseries, text=labelcolor},
    sublabel/.style={font=\tiny\sffamily, text=labelcolor}
]

% ============================================================================
% DOMINIO FONTE (Esquerda)
% ============================================================================
\node[sectionlabel] at (0, 3.2) {Domínio Fonte};
\node[sublabel] at (0, 2.8) {(Pré-treinamento)};

\node[datasetbox=inputcolor] (source_data) at (0, 1.5) {};
\node[font=\scriptsize\sffamily\bfseries, text=inputcolor!80!black] at (0, 1.8) {COCO};
\node[font=\tiny\sffamily, text=labelcolor] at (0, 1.2) {Milhões de imagens};

% Camadas fonte (compactas)
\node[layerfrozen] (src_conv1) at (0, -0.1) {Conv - Bordas};
\node[layerfrozen] (src_conv2) at (0, -0.7) {Conv - Texturas};
\node[layerfrozen] (src_conv3) at (0, -1.3) {Conv - Padrões};
\node[layerfinetune] (src_fc1) at (0, -1.9) {FC - Semântica};
\node[layernew] (src_fc2) at (0, -2.5) {FC - Classes};

\draw[arrow] (source_data.south) -- (src_conv1.north);
\draw[arrow] (src_conv1) -- (src_conv2);
\draw[arrow] (src_conv2) -- (src_conv3);
\draw[arrow] (src_conv3) -- (src_fc1);
\draw[arrow] (src_fc1) -- (src_fc2);

\node[sublabel] at (0, -3.1) {80 classes genéricas};

% ============================================================================
% PROCESSO DE TRANSFERENCIA (Centro)
% ============================================================================
\draw[transferarrow] (1.8, -1.0) -- node[above, font=\scriptsize\sffamily\bfseries, text=skipcolor] {Transfer} node[below, font=\tiny\sffamily, text=skipcolor] {Learning} (4.2, -1.0);

% ============================================================================
% DOMINIO ALVO (Direita)
% ============================================================================
\node[sectionlabel] at (6.0, 3.2) {Domínio Alvo};
\node[sublabel] at (6.0, 2.8) {(Fine-tuning)};

\node[datasetbox=headcolor] (target_data) at (6.0, 1.5) {};
\node[font=\scriptsize\sffamily\bfseries, text=headcolor!80!black] at (6.0, 1.8) {Dataset EPI};
\node[font=\tiny\sffamily, text=labelcolor] at (6.0, 1.2) {Classes específicas};

% Camadas alvo
\node[layerfrozen] (tgt_conv1) at (6.0, -0.1) {Conv - Bordas};
\node[layerfrozen] (tgt_conv2) at (6.0, -0.7) {Conv - Texturas};
\node[layerfinetune] (tgt_conv3) at (6.0, -1.3) {Conv - EPIs};
\node[layerfinetune] (tgt_fc1) at (6.0, -1.9) {FC - Atributos};
\node[layernew] (tgt_fc2) at (6.0, -2.5) {FC - Classes EPI};

\draw[arrow] (target_data.south) -- (tgt_conv1.north);
\draw[arrow] (tgt_conv1) -- (tgt_conv2);
\draw[arrow] (tgt_conv2) -- (tgt_conv3);
\draw[arrow] (tgt_conv3) -- (tgt_fc1);
\draw[arrow] (tgt_fc1) -- (tgt_fc2);

\node[sublabel] at (6.0, -3.1) {Capacete, Colete, Luvas...};

% ============================================================================
% LEGENDA COMPACTA
% ============================================================================
\begin{scope}[shift={(9.0, 0.5)}]
    \node[font=\footnotesize\sffamily\bfseries, text=labelcolor] at (0, 0.8) {Legenda:};

    \node[layerfrozen, minimum width=1.2cm] at (0, 0.3) {};
    \node[font=\tiny\sffamily, text=labelcolor, anchor=west] at (0.7, 0.3) {Congelada};

    \node[layerfinetune, minimum width=1.2cm] at (0, -0.2) {};
    \node[font=\tiny\sffamily, text=labelcolor, anchor=west] at (0.7, -0.2) {Ajustada};

    \node[layernew, minimum width=1.2cm] at (0, -0.7) {};
    \node[font=\tiny\sffamily, text=labelcolor, anchor=west] at (0.7, -0.7) {Nova};
\end{scope}

% Seta de conhecimento transferido
\draw[dashed, skipcolor, line width=0.6pt, -{Stealth[length=1.5mm]}]
    (0.8, -0.4) .. controls (2.5, 0.8) and (3.5, 0.8) .. (5.2, -0.4);
\node[font=\tiny\sffamily\itshape, text=skipcolor] at (3.0, 1.2) {Pesos transferidos};

\end{tikzpicture}

% ============================================================================

\begin{tikzpicture}[
    scale=0.85, transform shape,
    % Estilos dos blocos
    layerfrozen/.style={
        rectangle, rounded corners=2pt,
        minimum width=2.0cm, minimum height=0.45cm,
        text centered, font=\tiny\sffamily,
        fill=cspcolor, text=white,
        draw=cspcolor!80!black, line width=0.5pt
    },
    layerfinetune/.style={
        rectangle, rounded corners=2pt,
        minimum width=2.0cm, minimum height=0.45cm,
        text centered, font=\tiny\sffamily,
        fill=headcolor, text=white,
        draw=headcolor!80!black, line width=0.5pt
    },
    layernew/.style={
        rectangle, rounded corners=2pt,
        minimum width=2.0cm, minimum height=0.45cm,
        text centered, font=\tiny\sffamily,
        fill=outputcolor, text=white,
        draw=outputcolor!80!black, line width=0.5pt
    },
    datasetbox/.style={
        rectangle, rounded corners=4pt,
        minimum width=2.5cm, minimum height=1.4cm,
        text centered, font=\scriptsize\sffamily,
        fill=#1!10, draw=#1, line width=1pt
    },
    arrow/.style={-{Stealth[length=1.8mm, width=1.2mm]}, thick, arrowcolor},
    transferarrow/.style={-{Stealth[length=2.5mm, width=1.8mm]}, very thick, skipcolor},
    sectionlabel/.style={font=\small\sffamily\bfseries, text=labelcolor},
    sublabel/.style={font=\tiny\sffamily, text=labelcolor}
]

% ============================================================================
% DOMINIO FONTE (Esquerda)
% ============================================================================
\node[sectionlabel] at (0, 3.2) {Domínio Fonte};
\node[sublabel] at (0, 2.8) {(Pré-treinamento)};

\node[datasetbox=inputcolor] (source_data) at (0, 1.5) {};
\node[font=\scriptsize\sffamily\bfseries, text=inputcolor!80!black] at (0, 1.8) {COCO};
\node[font=\tiny\sffamily, text=labelcolor] at (0, 1.2) {Milhões de imagens};

% Camadas fonte (compactas)
\node[layerfrozen] (src_conv1) at (0, -0.1) {Conv - Bordas};
\node[layerfrozen] (src_conv2) at (0, -0.7) {Conv - Texturas};
\node[layerfrozen] (src_conv3) at (0, -1.3) {Conv - Padrões};
\node[layerfinetune] (src_fc1) at (0, -1.9) {FC - Semântica};
\node[layernew] (src_fc2) at (0, -2.5) {FC - Classes};

\draw[arrow] (source_data.south) -- (src_conv1.north);
\draw[arrow] (src_conv1) -- (src_conv2);
\draw[arrow] (src_conv2) -- (src_conv3);
\draw[arrow] (src_conv3) -- (src_fc1);
\draw[arrow] (src_fc1) -- (src_fc2);

\node[sublabel] at (0, -3.1) {80 classes genéricas};

% ============================================================================
% PROCESSO DE TRANSFERENCIA (Centro)
% ============================================================================
\draw[transferarrow] (1.8, -1.0) -- node[above, font=\scriptsize\sffamily\bfseries, text=skipcolor] {Transfer} node[below, font=\tiny\sffamily, text=skipcolor] {Learning} (4.2, -1.0);

% ============================================================================
% DOMINIO ALVO (Direita)
% ============================================================================
\node[sectionlabel] at (6.0, 3.2) {Domínio Alvo};
\node[sublabel] at (6.0, 2.8) {(Fine-tuning)};

\node[datasetbox=headcolor] (target_data) at (6.0, 1.5) {};
\node[font=\scriptsize\sffamily\bfseries, text=headcolor!80!black] at (6.0, 1.8) {Dataset EPI};
\node[font=\tiny\sffamily, text=labelcolor] at (6.0, 1.2) {Classes específicas};

% Camadas alvo
\node[layerfrozen] (tgt_conv1) at (6.0, -0.1) {Conv - Bordas};
\node[layerfrozen] (tgt_conv2) at (6.0, -0.7) {Conv - Texturas};
\node[layerfinetune] (tgt_conv3) at (6.0, -1.3) {Conv - EPIs};
\node[layerfinetune] (tgt_fc1) at (6.0, -1.9) {FC - Atributos};
\node[layernew] (tgt_fc2) at (6.0, -2.5) {FC - Classes EPI};

\draw[arrow] (target_data.south) -- (tgt_conv1.north);
\draw[arrow] (tgt_conv1) -- (tgt_conv2);
\draw[arrow] (tgt_conv2) -- (tgt_conv3);
\draw[arrow] (tgt_conv3) -- (tgt_fc1);
\draw[arrow] (tgt_fc1) -- (tgt_fc2);

\node[sublabel] at (6.0, -3.1) {Capacete, Colete, Luvas...};

% ============================================================================
% LEGENDA COMPACTA
% ============================================================================
\begin{scope}[shift={(9.0, 0.5)}]
    \node[font=\footnotesize\sffamily\bfseries, text=labelcolor] at (0, 0.8) {Legenda:};

    \node[layerfrozen, minimum width=1.2cm] at (0, 0.3) {};
    \node[font=\tiny\sffamily, text=labelcolor, anchor=west] at (0.7, 0.3) {Congelada};

    \node[layerfinetune, minimum width=1.2cm] at (0, -0.2) {};
    \node[font=\tiny\sffamily, text=labelcolor, anchor=west] at (0.7, -0.2) {Ajustada};

    \node[layernew, minimum width=1.2cm] at (0, -0.7) {};
    \node[font=\tiny\sffamily, text=labelcolor, anchor=west] at (0.7, -0.7) {Nova};
\end{scope}

% Seta de conhecimento transferido
\draw[dashed, skipcolor, line width=0.6pt, -{Stealth[length=1.5mm]}]
    (0.8, -0.4) .. controls (2.5, 0.8) and (3.5, 0.8) .. (5.2, -0.4);
\node[font=\tiny\sffamily\itshape, text=skipcolor] at (3.0, 1.2) {Pesos transferidos};

\end{tikzpicture}
}
  \end{center}
  \vfill

  \begin{beamercolorbox}[wd=\textwidth, rounded=true, shadow=false]{block body}
    \centering
    {\footnotesize\color{stagecolor}%
      Essencial para \textit{datasets} pequenos: camadas iniciais (bordas, texturas) são \textbf{reutilizadas}, apenas camadas finais são ajustadas}
  \end{beamercolorbox}
\end{frame}

% --- Slide: Non-Maximum Suppression ---
\begin{frame}{Non-Maximum Suppression (NMS)}
  \vspace{-2mm}
  {\small\color{sublabelcolor}%
    Supressão de detecções redundantes: mantém apenas a predição de \textbf{maior confiança} por objeto}

  \vfill
  \begin{center}
    \resizebox{!}{0.58\textheight}{% ============================================================================
% Diagrama do Algoritmo Non-Maximum Suppression (NMS) - Simplificado
% Para uso em dissertação de mestrado - PPGC/UFRGS
% Uso: % ============================================================================
% Diagrama do Algoritmo Non-Maximum Suppression (NMS) - Simplificado
% Para uso em dissertação de mestrado - PPGC/UFRGS
% Uso: % ============================================================================
% Diagrama do Algoritmo Non-Maximum Suppression (NMS) - Simplificado
% Para uso em dissertação de mestrado - PPGC/UFRGS
% Uso: \input{figuras/diagrama_nms.tex}
% ============================================================================

\begin{tikzpicture}[
    node distance=0.6cm and 1.0cm,
    scale=0.9, transform shape,
    font=\sffamily,
    % Estilos
    sectionbox/.style={
        rectangle,
        rounded corners=4pt,
        draw=labelcolor,
        fill=bgbox,
        minimum width=3.8cm,
        minimum height=4.5cm,
        line width=0.6pt,
    },
    processbox/.style={
        rectangle,
        rounded corners=6pt,
        fill=cspcolor,
        text=white,
        font=\normalsize\sffamily\bfseries,
        minimum width=2.4cm,
        minimum height=1.0cm,
        align=center,
    },
    sectiontitle/.style={
        font=\normalsize\sffamily\bfseries,
        text=labelcolor,
    },
    processarrow/.style={
        -{Stealth[length=3mm, width=2mm]},
        line width=1.5pt,
        arrowcolor,
    },
]

% Cores para bounding boxes
\definecolor{bbox1}{RGB}{39, 174, 96}
\definecolor{bbox2}{RGB}{41, 128, 185}
\definecolor{bbox3}{RGB}{230, 126, 34}
\definecolor{bbox4}{RGB}{192, 57, 43}

% ============================================================================
% ANTES DO NMS
% ============================================================================
\node[sectionbox] (beforebox) at (0, 0) {};
\node[sectiontitle, above=0.05cm of beforebox.north] {Antes do NMS};

% Silhueta simplificada
\begin{scope}[shift={(0, -0.2)}]
    \fill[labelcolor!50] (0, 1.1) circle (0.28);
    \fill[labelcolor!50, rounded corners=2pt] (-0.32, -0.7) rectangle (0.32, 0.75);
    \fill[labelcolor!50, rounded corners=2pt] (-0.28, -0.7) rectangle (-0.08, -1.4);
    \fill[labelcolor!50, rounded corners=2pt] (0.08, -0.7) rectangle (0.28, -1.4);
\end{scope}

% Bounding boxes sobrepostas
\draw[bbox4, line width=1pt, dashed] (-1.1, -1.8) rectangle (1.0, 1.6);
\node[font=\scriptsize\bfseries, fill=bbox4, text=white, rounded corners=1pt, inner sep=1pt] at (1.0, 1.6) {0.71};

\draw[bbox3, line width=1.2pt, densely dashed] (-0.95, -1.65) rectangle (0.85, 1.5);
\node[font=\scriptsize\bfseries, fill=bbox3, text=white, rounded corners=1pt, inner sep=1pt] at (0.85, 1.5) {0.78};

\draw[bbox2, line width=1.3pt, densely dotted] (-0.8, -1.55) rectangle (0.7, 1.45);
\node[font=\scriptsize\bfseries, fill=bbox2, text=white, rounded corners=1pt, inner sep=1pt] at (0.7, 1.45) {0.85};

\draw[bbox1, line width=1.5pt] (-0.7, -1.5) rectangle (0.6, 1.4);
\node[font=\scriptsize\bfseries, fill=bbox1, text=white, rounded corners=1pt, inner sep=1pt] at (0.6, 1.4) {0.92};

\node[font=\scriptsize\sffamily\itshape, text=labelcolor] at (0, -2.2) {Múltiplas detecções};

% ============================================================================
% PROCESSO NMS
% ============================================================================
\node[processbox] (nmsprocess) at (4.5, 0) {NMS\\IoU $>$ 0.5};
\draw[processarrow] (beforebox.east) -- (nmsprocess.west);
\draw[processarrow] (nmsprocess.east) -- ++(1.5, 0);

\node[font=\scriptsize\sffamily, text=labelcolor, align=center] at (4.5, -1.0) {Suprime caixas\\redundantes};

% ============================================================================
% DEPOIS DO NMS
% ============================================================================
\node[sectionbox] (afterbox) at (9.0, 0) {};
\node[sectiontitle, above=0.05cm of afterbox.north] {Após o NMS};

% Silhueta
\begin{scope}[shift={(9.0, -0.2)}]
    \fill[labelcolor!50] (0, 1.1) circle (0.28);
    \fill[labelcolor!50, rounded corners=2pt] (-0.32, -0.7) rectangle (0.32, 0.75);
    \fill[labelcolor!50, rounded corners=2pt] (-0.28, -0.7) rectangle (-0.08, -1.4);
    \fill[labelcolor!50, rounded corners=2pt] (0.08, -0.7) rectangle (0.28, -1.4);
\end{scope}

% Apenas melhor box
\draw[bbox1, line width=1.8pt] (8.3, -1.5) rectangle (9.6, 1.4);
\node[font=\scriptsize\bfseries, fill=bbox1, text=white, rounded corners=1pt, inner sep=1pt] at (9.6, 1.4) {0.92};

\node[font=\scriptsize\sffamily\itshape, text=labelcolor] at (9.0, -2.2) {Melhor detecção};

% ============================================================================
% LEGENDA COMPACTA
% ============================================================================
\node[font=\small\sffamily\bfseries, text=labelcolor] at (4.5, -2.8) {Etapas:};
\node[font=\scriptsize\sffamily, text=labelcolor, align=left] at (4.5, -3.5) {
    1. Ordenar por confiança\\
    2. Manter a melhor caixa\\
    3. Suprimir IoU $>$ limiar\\
    4. Repetir para restantes
};

\end{tikzpicture}

% ============================================================================

\begin{tikzpicture}[
    node distance=0.6cm and 1.0cm,
    scale=0.9, transform shape,
    font=\sffamily,
    % Estilos
    sectionbox/.style={
        rectangle,
        rounded corners=4pt,
        draw=labelcolor,
        fill=bgbox,
        minimum width=3.8cm,
        minimum height=4.5cm,
        line width=0.6pt,
    },
    processbox/.style={
        rectangle,
        rounded corners=6pt,
        fill=cspcolor,
        text=white,
        font=\normalsize\sffamily\bfseries,
        minimum width=2.4cm,
        minimum height=1.0cm,
        align=center,
    },
    sectiontitle/.style={
        font=\normalsize\sffamily\bfseries,
        text=labelcolor,
    },
    processarrow/.style={
        -{Stealth[length=3mm, width=2mm]},
        line width=1.5pt,
        arrowcolor,
    },
]

% Cores para bounding boxes
\definecolor{bbox1}{RGB}{39, 174, 96}
\definecolor{bbox2}{RGB}{41, 128, 185}
\definecolor{bbox3}{RGB}{230, 126, 34}
\definecolor{bbox4}{RGB}{192, 57, 43}

% ============================================================================
% ANTES DO NMS
% ============================================================================
\node[sectionbox] (beforebox) at (0, 0) {};
\node[sectiontitle, above=0.05cm of beforebox.north] {Antes do NMS};

% Silhueta simplificada
\begin{scope}[shift={(0, -0.2)}]
    \fill[labelcolor!50] (0, 1.1) circle (0.28);
    \fill[labelcolor!50, rounded corners=2pt] (-0.32, -0.7) rectangle (0.32, 0.75);
    \fill[labelcolor!50, rounded corners=2pt] (-0.28, -0.7) rectangle (-0.08, -1.4);
    \fill[labelcolor!50, rounded corners=2pt] (0.08, -0.7) rectangle (0.28, -1.4);
\end{scope}

% Bounding boxes sobrepostas
\draw[bbox4, line width=1pt, dashed] (-1.1, -1.8) rectangle (1.0, 1.6);
\node[font=\scriptsize\bfseries, fill=bbox4, text=white, rounded corners=1pt, inner sep=1pt] at (1.0, 1.6) {0.71};

\draw[bbox3, line width=1.2pt, densely dashed] (-0.95, -1.65) rectangle (0.85, 1.5);
\node[font=\scriptsize\bfseries, fill=bbox3, text=white, rounded corners=1pt, inner sep=1pt] at (0.85, 1.5) {0.78};

\draw[bbox2, line width=1.3pt, densely dotted] (-0.8, -1.55) rectangle (0.7, 1.45);
\node[font=\scriptsize\bfseries, fill=bbox2, text=white, rounded corners=1pt, inner sep=1pt] at (0.7, 1.45) {0.85};

\draw[bbox1, line width=1.5pt] (-0.7, -1.5) rectangle (0.6, 1.4);
\node[font=\scriptsize\bfseries, fill=bbox1, text=white, rounded corners=1pt, inner sep=1pt] at (0.6, 1.4) {0.92};

\node[font=\scriptsize\sffamily\itshape, text=labelcolor] at (0, -2.2) {Múltiplas detecções};

% ============================================================================
% PROCESSO NMS
% ============================================================================
\node[processbox] (nmsprocess) at (4.5, 0) {NMS\\IoU $>$ 0.5};
\draw[processarrow] (beforebox.east) -- (nmsprocess.west);
\draw[processarrow] (nmsprocess.east) -- ++(1.5, 0);

\node[font=\scriptsize\sffamily, text=labelcolor, align=center] at (4.5, -1.0) {Suprime caixas\\redundantes};

% ============================================================================
% DEPOIS DO NMS
% ============================================================================
\node[sectionbox] (afterbox) at (9.0, 0) {};
\node[sectiontitle, above=0.05cm of afterbox.north] {Após o NMS};

% Silhueta
\begin{scope}[shift={(9.0, -0.2)}]
    \fill[labelcolor!50] (0, 1.1) circle (0.28);
    \fill[labelcolor!50, rounded corners=2pt] (-0.32, -0.7) rectangle (0.32, 0.75);
    \fill[labelcolor!50, rounded corners=2pt] (-0.28, -0.7) rectangle (-0.08, -1.4);
    \fill[labelcolor!50, rounded corners=2pt] (0.08, -0.7) rectangle (0.28, -1.4);
\end{scope}

% Apenas melhor box
\draw[bbox1, line width=1.8pt] (8.3, -1.5) rectangle (9.6, 1.4);
\node[font=\scriptsize\bfseries, fill=bbox1, text=white, rounded corners=1pt, inner sep=1pt] at (9.6, 1.4) {0.92};

\node[font=\scriptsize\sffamily\itshape, text=labelcolor] at (9.0, -2.2) {Melhor detecção};

% ============================================================================
% LEGENDA COMPACTA
% ============================================================================
\node[font=\small\sffamily\bfseries, text=labelcolor] at (4.5, -2.8) {Etapas:};
\node[font=\scriptsize\sffamily, text=labelcolor, align=left] at (4.5, -3.5) {
    1. Ordenar por confiança\\
    2. Manter a melhor caixa\\
    3. Suprimir IoU $>$ limiar\\
    4. Repetir para restantes
};

\end{tikzpicture}

% ============================================================================

\begin{tikzpicture}[
    node distance=0.6cm and 1.0cm,
    scale=0.9, transform shape,
    font=\sffamily,
    % Estilos
    sectionbox/.style={
        rectangle,
        rounded corners=4pt,
        draw=labelcolor,
        fill=bgbox,
        minimum width=3.8cm,
        minimum height=4.5cm,
        line width=0.6pt,
    },
    processbox/.style={
        rectangle,
        rounded corners=6pt,
        fill=cspcolor,
        text=white,
        font=\normalsize\sffamily\bfseries,
        minimum width=2.4cm,
        minimum height=1.0cm,
        align=center,
    },
    sectiontitle/.style={
        font=\normalsize\sffamily\bfseries,
        text=labelcolor,
    },
    processarrow/.style={
        -{Stealth[length=3mm, width=2mm]},
        line width=1.5pt,
        arrowcolor,
    },
]

% Cores para bounding boxes
\definecolor{bbox1}{RGB}{39, 174, 96}
\definecolor{bbox2}{RGB}{41, 128, 185}
\definecolor{bbox3}{RGB}{230, 126, 34}
\definecolor{bbox4}{RGB}{192, 57, 43}

% ============================================================================
% ANTES DO NMS
% ============================================================================
\node[sectionbox] (beforebox) at (0, 0) {};
\node[sectiontitle, above=0.05cm of beforebox.north] {Antes do NMS};

% Silhueta simplificada
\begin{scope}[shift={(0, -0.2)}]
    \fill[labelcolor!50] (0, 1.1) circle (0.28);
    \fill[labelcolor!50, rounded corners=2pt] (-0.32, -0.7) rectangle (0.32, 0.75);
    \fill[labelcolor!50, rounded corners=2pt] (-0.28, -0.7) rectangle (-0.08, -1.4);
    \fill[labelcolor!50, rounded corners=2pt] (0.08, -0.7) rectangle (0.28, -1.4);
\end{scope}

% Bounding boxes sobrepostas
\draw[bbox4, line width=1pt, dashed] (-1.1, -1.8) rectangle (1.0, 1.6);
\node[font=\scriptsize\bfseries, fill=bbox4, text=white, rounded corners=1pt, inner sep=1pt] at (1.0, 1.6) {0.71};

\draw[bbox3, line width=1.2pt, densely dashed] (-0.95, -1.65) rectangle (0.85, 1.5);
\node[font=\scriptsize\bfseries, fill=bbox3, text=white, rounded corners=1pt, inner sep=1pt] at (0.85, 1.5) {0.78};

\draw[bbox2, line width=1.3pt, densely dotted] (-0.8, -1.55) rectangle (0.7, 1.45);
\node[font=\scriptsize\bfseries, fill=bbox2, text=white, rounded corners=1pt, inner sep=1pt] at (0.7, 1.45) {0.85};

\draw[bbox1, line width=1.5pt] (-0.7, -1.5) rectangle (0.6, 1.4);
\node[font=\scriptsize\bfseries, fill=bbox1, text=white, rounded corners=1pt, inner sep=1pt] at (0.6, 1.4) {0.92};

\node[font=\scriptsize\sffamily\itshape, text=labelcolor] at (0, -2.2) {Múltiplas detecções};

% ============================================================================
% PROCESSO NMS
% ============================================================================
\node[processbox] (nmsprocess) at (4.5, 0) {NMS\\IoU $>$ 0.5};
\draw[processarrow] (beforebox.east) -- (nmsprocess.west);
\draw[processarrow] (nmsprocess.east) -- ++(1.5, 0);

\node[font=\scriptsize\sffamily, text=labelcolor, align=center] at (4.5, -1.0) {Suprime caixas\\redundantes};

% ============================================================================
% DEPOIS DO NMS
% ============================================================================
\node[sectionbox] (afterbox) at (9.0, 0) {};
\node[sectiontitle, above=0.05cm of afterbox.north] {Após o NMS};

% Silhueta
\begin{scope}[shift={(9.0, -0.2)}]
    \fill[labelcolor!50] (0, 1.1) circle (0.28);
    \fill[labelcolor!50, rounded corners=2pt] (-0.32, -0.7) rectangle (0.32, 0.75);
    \fill[labelcolor!50, rounded corners=2pt] (-0.28, -0.7) rectangle (-0.08, -1.4);
    \fill[labelcolor!50, rounded corners=2pt] (0.08, -0.7) rectangle (0.28, -1.4);
\end{scope}

% Apenas melhor box
\draw[bbox1, line width=1.8pt] (8.3, -1.5) rectangle (9.6, 1.4);
\node[font=\scriptsize\bfseries, fill=bbox1, text=white, rounded corners=1pt, inner sep=1pt] at (9.6, 1.4) {0.92};

\node[font=\scriptsize\sffamily\itshape, text=labelcolor] at (9.0, -2.2) {Melhor detecção};

% ============================================================================
% LEGENDA COMPACTA
% ============================================================================
\node[font=\small\sffamily\bfseries, text=labelcolor] at (4.5, -2.8) {Etapas:};
\node[font=\scriptsize\sffamily, text=labelcolor, align=left] at (4.5, -3.5) {
    1. Ordenar por confiança\\
    2. Manter a melhor caixa\\
    3. Suprimir IoU $>$ limiar\\
    4. Repetir para restantes
};

\end{tikzpicture}
}
  \end{center}
  \vfill

  \begin{beamercolorbox}[wd=\textwidth, rounded=true, shadow=false]{block body}
    \centering
    {\footnotesize\color{stagecolor}%
      Limiar \textbf{IoU $>$ 0,5}: caixas com sobreposição acima desse valor são suprimidas --- etapa final do \textit{pipeline} YOLO}
  \end{beamercolorbox}
\end{frame}

% --- Slide: Jetson AGX Xavier ---
\begin{frame}{NVIDIA Jetson AGX Xavier}
  \vspace{-2mm}
  {\small\color{sublabelcolor}%
    Plataforma de computação de borda para inferência de IA em tempo real}

  \vspace{2mm}
  \begin{columns}[T]
    \begin{column}{0.58\textwidth}
      \begin{center}
        \resizebox{\textwidth}{!}{% ============================================================================
% Diagrama de Arquitetura de Hardware - NVIDIA Jetson AGX Xavier
% Para uso em dissertacao de mestrado PPGC/UFRGS
% Uso: % ============================================================================
% Diagrama de Arquitetura de Hardware - NVIDIA Jetson AGX Xavier
% Para uso em dissertacao de mestrado PPGC/UFRGS
% Uso: % ============================================================================
% Diagrama de Arquitetura de Hardware - NVIDIA Jetson AGX Xavier
% Para uso em dissertacao de mestrado PPGC/UFRGS
% Uso: \input{figuras/diagrama_jetson.tex}
% ============================================================================

\begin{tikzpicture}[
    % Estilos para os blocos principais
    mainblock/.style={
        rectangle,
        rounded corners=3pt,
        draw=#1!80!black,
        fill=#1!20,
        line width=1pt,
        minimum width=4.2cm,
        minimum height=2.2cm,
        align=center,
        font=\sffamily\small
    },
    % Estilo para blocos menores
    smallblock/.style={
        rectangle,
        rounded corners=2pt,
        draw=#1!80!black,
        fill=#1!15,
        line width=0.8pt,
        minimum width=2.8cm,
        minimum height=1.4cm,
        align=center,
        font=\sffamily\footnotesize
    },
    % Estilo para a memoria compartilhada
    memblock/.style={
        rectangle,
        rounded corners=3pt,
        draw=memcolor!80!black,
        fill=memcolor!15,
        line width=1.2pt,
        minimum width=10cm,
        minimum height=1.2cm,
        align=center,
        font=\sffamily\small\bfseries
    },
    % Estilo para interfaces I/O
    ioblock/.style={
        rectangle,
        rounded corners=2pt,
        draw=iocolor!80!black,
        fill=iocolor!10,
        line width=0.6pt,
        minimum width=1.8cm,
        minimum height=0.9cm,
        align=center,
        font=\sffamily\scriptsize
    },
    % Estilo para setas de conexao
    busarrow/.style={
        -{Stealth[length=3mm, width=2mm]},
        line width=1.5pt,
        draw=buscolor
    },
    doublearrow/.style={
        {Stealth[length=2.5mm, width=2mm]}-{Stealth[length=2.5mm, width=2mm]},
        line width=1.2pt,
        draw=buscolor
    },
    % Estilo para labels
    labelstyle/.style={
        font=\sffamily\scriptsize,
        text=black!70
    }
]

% ============================================================================
% Bloco de Memoria Compartilhada (topo)
% ============================================================================
\node[memblock] (mem) at (0, 4.5) {
    \textbf{Memoria Unificada LPDDR4x}\\
    32 GB | 136.5 GB/s
};

% ============================================================================
% CPU - Lado esquerdo
% ============================================================================
\node[mainblock=cpucolor] (cpu) at (-3.5, 1.5) {
    \textbf{CPU}\\[2pt]
    8 nucleos ARM Carmel\\
    64-bit v8.2\\
    2.26 GHz
};

% ============================================================================
% GPU - Lado direito
% ============================================================================
\node[mainblock=gpucolor] (gpu) at (3.5, 1.5) {
    \textbf{GPU Volta}\\[2pt]
    512 CUDA Cores\\
    64 Tensor Cores\\
    32 TOPS (INT8)
};

% ============================================================================
% DLA - Deep Learning Accelerators
% ============================================================================
\node[smallblock=dlacolor] (dla) at (-3.5, -1.5) {
    \textbf{DLA}\\[2pt]
    2x Deep Learning\\
    Accelerators
};

% ============================================================================
% Video Engine
% ============================================================================
\node[smallblock=videocolor] (video) at (0, -1.5) {
    \textbf{Motor de Video}\\[2pt]
    NVENC/NVDEC\\
    H.264/H.265
};

% ============================================================================
% Vision Accelerator
% ============================================================================
\node[smallblock=dlacolor] (vision) at (3.5, -1.5) {
    \textbf{Vision Accelerator}\\[2pt]
    PVA\\
    7-Way VLIW
};

% ============================================================================
% Interfaces de I/O (parte inferior)
% ============================================================================
\node[ioblock] (csi) at (-4.5, -4) {
    \textbf{CSI}\\
    16 Cameras
};

\node[ioblock] (usb) at (-2.25, -4) {
    \textbf{USB}\\
    3.0/2.0
};

\node[ioblock] (pcie) at (0, -4) {
    \textbf{PCIe}\\
    Gen4 x8
};

\node[ioblock] (eth) at (2.25, -4) {
    \textbf{Ethernet}\\
    10 GbE
};

\node[ioblock] (gpio) at (4.5, -4) {
    \textbf{GPIO}\\
    40 pinos
};

% ============================================================================
% Conexoes - Memoria para CPU e GPU
% ============================================================================
\draw[doublearrow] (mem.south -| cpu.north) -- (cpu.north);
\draw[doublearrow] (mem.south -| gpu.north) -- (gpu.north);

% ============================================================================
% Conexoes - CPU/GPU para blocos inferiores via barramento
% ============================================================================
% Linha de barramento horizontal
\draw[line width=2pt, draw=buscolor!60] (-5, -0.3) -- (5, -0.3);

% Conexoes verticais para o barramento
\draw[line width=1.2pt, draw=buscolor!60] (cpu.south) -- (cpu.south |- 0,-0.3);
\draw[line width=1.2pt, draw=buscolor!60] (gpu.south) -- (gpu.south |- 0,-0.3);
\draw[line width=1.2pt, draw=buscolor!60] (dla.north) -- (dla.north |- 0,-0.3);
\draw[line width=1.2pt, draw=buscolor!60] (video.north) -- (video.north |- 0,-0.3);
\draw[line width=1.2pt, draw=buscolor!60] (vision.north) -- (vision.north |- 0,-0.3);

% ============================================================================
% Conexoes - Blocos para I/O
% ============================================================================
% Linha de barramento I/O
\draw[line width=1.5pt, draw=iocolor!60] (-5, -2.8) -- (5, -2.8);

% Conexoes para I/O
\draw[line width=0.8pt, draw=iocolor!60] (dla.south) -- (dla.south |- 0,-2.8);
\draw[line width=0.8pt, draw=iocolor!60] (video.south) -- (video.south |- 0,-2.8);
\draw[line width=0.8pt, draw=iocolor!60] (vision.south) -- (vision.south |- 0,-2.8);

\draw[line width=0.8pt, draw=iocolor!60] (csi.north) -- (csi.north |- 0,-2.8);
\draw[line width=0.8pt, draw=iocolor!60] (usb.north) -- (usb.north |- 0,-2.8);
\draw[line width=0.8pt, draw=iocolor!60] (pcie.north) -- (pcie.north |- 0,-2.8);
\draw[line width=0.8pt, draw=iocolor!60] (eth.north) -- (eth.north |- 0,-2.8);
\draw[line width=0.8pt, draw=iocolor!60] (gpio.north) -- (gpio.north |- 0,-2.8);

% ============================================================================
% Labels para barramentos
% ============================================================================
\node[labelstyle, anchor=east] at (-5.2, -0.3) {Barramento};
\node[labelstyle, anchor=east] at (-5.2, -0.6) {de Sistema};

\node[labelstyle, anchor=east] at (-5.2, -2.8) {Barramento};
\node[labelstyle, anchor=east] at (-5.2, -3.1) {de I/O};

% ============================================================================
% Titulo do diagrama
% ============================================================================
\node[font=\sffamily\large\bfseries, anchor=south] at (0, 5.3) {
    Arquitetura NVIDIA Jetson AGX Xavier
};

% ============================================================================
% Legenda de potencia
% ============================================================================
\node[font=\sffamily\scriptsize, text=black!60, anchor=north west] at (-5.5, -4.8) {
    Consumo: 10W -- 30W (configuravel)
};

\end{tikzpicture}

% ============================================================================

\begin{tikzpicture}[
    % Estilos para os blocos principais
    mainblock/.style={
        rectangle,
        rounded corners=3pt,
        draw=#1!80!black,
        fill=#1!20,
        line width=1pt,
        minimum width=4.2cm,
        minimum height=2.2cm,
        align=center,
        font=\sffamily\small
    },
    % Estilo para blocos menores
    smallblock/.style={
        rectangle,
        rounded corners=2pt,
        draw=#1!80!black,
        fill=#1!15,
        line width=0.8pt,
        minimum width=2.8cm,
        minimum height=1.4cm,
        align=center,
        font=\sffamily\footnotesize
    },
    % Estilo para a memoria compartilhada
    memblock/.style={
        rectangle,
        rounded corners=3pt,
        draw=memcolor!80!black,
        fill=memcolor!15,
        line width=1.2pt,
        minimum width=10cm,
        minimum height=1.2cm,
        align=center,
        font=\sffamily\small\bfseries
    },
    % Estilo para interfaces I/O
    ioblock/.style={
        rectangle,
        rounded corners=2pt,
        draw=iocolor!80!black,
        fill=iocolor!10,
        line width=0.6pt,
        minimum width=1.8cm,
        minimum height=0.9cm,
        align=center,
        font=\sffamily\scriptsize
    },
    % Estilo para setas de conexao
    busarrow/.style={
        -{Stealth[length=3mm, width=2mm]},
        line width=1.5pt,
        draw=buscolor
    },
    doublearrow/.style={
        {Stealth[length=2.5mm, width=2mm]}-{Stealth[length=2.5mm, width=2mm]},
        line width=1.2pt,
        draw=buscolor
    },
    % Estilo para labels
    labelstyle/.style={
        font=\sffamily\scriptsize,
        text=black!70
    }
]

% ============================================================================
% Bloco de Memoria Compartilhada (topo)
% ============================================================================
\node[memblock] (mem) at (0, 4.5) {
    \textbf{Memoria Unificada LPDDR4x}\\
    32 GB | 136.5 GB/s
};

% ============================================================================
% CPU - Lado esquerdo
% ============================================================================
\node[mainblock=cpucolor] (cpu) at (-3.5, 1.5) {
    \textbf{CPU}\\[2pt]
    8 nucleos ARM Carmel\\
    64-bit v8.2\\
    2.26 GHz
};

% ============================================================================
% GPU - Lado direito
% ============================================================================
\node[mainblock=gpucolor] (gpu) at (3.5, 1.5) {
    \textbf{GPU Volta}\\[2pt]
    512 CUDA Cores\\
    64 Tensor Cores\\
    32 TOPS (INT8)
};

% ============================================================================
% DLA - Deep Learning Accelerators
% ============================================================================
\node[smallblock=dlacolor] (dla) at (-3.5, -1.5) {
    \textbf{DLA}\\[2pt]
    2x Deep Learning\\
    Accelerators
};

% ============================================================================
% Video Engine
% ============================================================================
\node[smallblock=videocolor] (video) at (0, -1.5) {
    \textbf{Motor de Video}\\[2pt]
    NVENC/NVDEC\\
    H.264/H.265
};

% ============================================================================
% Vision Accelerator
% ============================================================================
\node[smallblock=dlacolor] (vision) at (3.5, -1.5) {
    \textbf{Vision Accelerator}\\[2pt]
    PVA\\
    7-Way VLIW
};

% ============================================================================
% Interfaces de I/O (parte inferior)
% ============================================================================
\node[ioblock] (csi) at (-4.5, -4) {
    \textbf{CSI}\\
    16 Cameras
};

\node[ioblock] (usb) at (-2.25, -4) {
    \textbf{USB}\\
    3.0/2.0
};

\node[ioblock] (pcie) at (0, -4) {
    \textbf{PCIe}\\
    Gen4 x8
};

\node[ioblock] (eth) at (2.25, -4) {
    \textbf{Ethernet}\\
    10 GbE
};

\node[ioblock] (gpio) at (4.5, -4) {
    \textbf{GPIO}\\
    40 pinos
};

% ============================================================================
% Conexoes - Memoria para CPU e GPU
% ============================================================================
\draw[doublearrow] (mem.south -| cpu.north) -- (cpu.north);
\draw[doublearrow] (mem.south -| gpu.north) -- (gpu.north);

% ============================================================================
% Conexoes - CPU/GPU para blocos inferiores via barramento
% ============================================================================
% Linha de barramento horizontal
\draw[line width=2pt, draw=buscolor!60] (-5, -0.3) -- (5, -0.3);

% Conexoes verticais para o barramento
\draw[line width=1.2pt, draw=buscolor!60] (cpu.south) -- (cpu.south |- 0,-0.3);
\draw[line width=1.2pt, draw=buscolor!60] (gpu.south) -- (gpu.south |- 0,-0.3);
\draw[line width=1.2pt, draw=buscolor!60] (dla.north) -- (dla.north |- 0,-0.3);
\draw[line width=1.2pt, draw=buscolor!60] (video.north) -- (video.north |- 0,-0.3);
\draw[line width=1.2pt, draw=buscolor!60] (vision.north) -- (vision.north |- 0,-0.3);

% ============================================================================
% Conexoes - Blocos para I/O
% ============================================================================
% Linha de barramento I/O
\draw[line width=1.5pt, draw=iocolor!60] (-5, -2.8) -- (5, -2.8);

% Conexoes para I/O
\draw[line width=0.8pt, draw=iocolor!60] (dla.south) -- (dla.south |- 0,-2.8);
\draw[line width=0.8pt, draw=iocolor!60] (video.south) -- (video.south |- 0,-2.8);
\draw[line width=0.8pt, draw=iocolor!60] (vision.south) -- (vision.south |- 0,-2.8);

\draw[line width=0.8pt, draw=iocolor!60] (csi.north) -- (csi.north |- 0,-2.8);
\draw[line width=0.8pt, draw=iocolor!60] (usb.north) -- (usb.north |- 0,-2.8);
\draw[line width=0.8pt, draw=iocolor!60] (pcie.north) -- (pcie.north |- 0,-2.8);
\draw[line width=0.8pt, draw=iocolor!60] (eth.north) -- (eth.north |- 0,-2.8);
\draw[line width=0.8pt, draw=iocolor!60] (gpio.north) -- (gpio.north |- 0,-2.8);

% ============================================================================
% Labels para barramentos
% ============================================================================
\node[labelstyle, anchor=east] at (-5.2, -0.3) {Barramento};
\node[labelstyle, anchor=east] at (-5.2, -0.6) {de Sistema};

\node[labelstyle, anchor=east] at (-5.2, -2.8) {Barramento};
\node[labelstyle, anchor=east] at (-5.2, -3.1) {de I/O};

% ============================================================================
% Titulo do diagrama
% ============================================================================
\node[font=\sffamily\large\bfseries, anchor=south] at (0, 5.3) {
    Arquitetura NVIDIA Jetson AGX Xavier
};

% ============================================================================
% Legenda de potencia
% ============================================================================
\node[font=\sffamily\scriptsize, text=black!60, anchor=north west] at (-5.5, -4.8) {
    Consumo: 10W -- 30W (configuravel)
};

\end{tikzpicture}

% ============================================================================

\begin{tikzpicture}[
    % Estilos para os blocos principais
    mainblock/.style={
        rectangle,
        rounded corners=3pt,
        draw=#1!80!black,
        fill=#1!20,
        line width=1pt,
        minimum width=4.2cm,
        minimum height=2.2cm,
        align=center,
        font=\sffamily\small
    },
    % Estilo para blocos menores
    smallblock/.style={
        rectangle,
        rounded corners=2pt,
        draw=#1!80!black,
        fill=#1!15,
        line width=0.8pt,
        minimum width=2.8cm,
        minimum height=1.4cm,
        align=center,
        font=\sffamily\footnotesize
    },
    % Estilo para a memoria compartilhada
    memblock/.style={
        rectangle,
        rounded corners=3pt,
        draw=memcolor!80!black,
        fill=memcolor!15,
        line width=1.2pt,
        minimum width=10cm,
        minimum height=1.2cm,
        align=center,
        font=\sffamily\small\bfseries
    },
    % Estilo para interfaces I/O
    ioblock/.style={
        rectangle,
        rounded corners=2pt,
        draw=iocolor!80!black,
        fill=iocolor!10,
        line width=0.6pt,
        minimum width=1.8cm,
        minimum height=0.9cm,
        align=center,
        font=\sffamily\scriptsize
    },
    % Estilo para setas de conexao
    busarrow/.style={
        -{Stealth[length=3mm, width=2mm]},
        line width=1.5pt,
        draw=buscolor
    },
    doublearrow/.style={
        {Stealth[length=2.5mm, width=2mm]}-{Stealth[length=2.5mm, width=2mm]},
        line width=1.2pt,
        draw=buscolor
    },
    % Estilo para labels
    labelstyle/.style={
        font=\sffamily\scriptsize,
        text=black!70
    }
]

% ============================================================================
% Bloco de Memoria Compartilhada (topo)
% ============================================================================
\node[memblock] (mem) at (0, 4.5) {
    \textbf{Memoria Unificada LPDDR4x}\\
    32 GB | 136.5 GB/s
};

% ============================================================================
% CPU - Lado esquerdo
% ============================================================================
\node[mainblock=cpucolor] (cpu) at (-3.5, 1.5) {
    \textbf{CPU}\\[2pt]
    8 nucleos ARM Carmel\\
    64-bit v8.2\\
    2.26 GHz
};

% ============================================================================
% GPU - Lado direito
% ============================================================================
\node[mainblock=gpucolor] (gpu) at (3.5, 1.5) {
    \textbf{GPU Volta}\\[2pt]
    512 CUDA Cores\\
    64 Tensor Cores\\
    32 TOPS (INT8)
};

% ============================================================================
% DLA - Deep Learning Accelerators
% ============================================================================
\node[smallblock=dlacolor] (dla) at (-3.5, -1.5) {
    \textbf{DLA}\\[2pt]
    2x Deep Learning\\
    Accelerators
};

% ============================================================================
% Video Engine
% ============================================================================
\node[smallblock=videocolor] (video) at (0, -1.5) {
    \textbf{Motor de Video}\\[2pt]
    NVENC/NVDEC\\
    H.264/H.265
};

% ============================================================================
% Vision Accelerator
% ============================================================================
\node[smallblock=dlacolor] (vision) at (3.5, -1.5) {
    \textbf{Vision Accelerator}\\[2pt]
    PVA\\
    7-Way VLIW
};

% ============================================================================
% Interfaces de I/O (parte inferior)
% ============================================================================
\node[ioblock] (csi) at (-4.5, -4) {
    \textbf{CSI}\\
    16 Cameras
};

\node[ioblock] (usb) at (-2.25, -4) {
    \textbf{USB}\\
    3.0/2.0
};

\node[ioblock] (pcie) at (0, -4) {
    \textbf{PCIe}\\
    Gen4 x8
};

\node[ioblock] (eth) at (2.25, -4) {
    \textbf{Ethernet}\\
    10 GbE
};

\node[ioblock] (gpio) at (4.5, -4) {
    \textbf{GPIO}\\
    40 pinos
};

% ============================================================================
% Conexoes - Memoria para CPU e GPU
% ============================================================================
\draw[doublearrow] (mem.south -| cpu.north) -- (cpu.north);
\draw[doublearrow] (mem.south -| gpu.north) -- (gpu.north);

% ============================================================================
% Conexoes - CPU/GPU para blocos inferiores via barramento
% ============================================================================
% Linha de barramento horizontal
\draw[line width=2pt, draw=buscolor!60] (-5, -0.3) -- (5, -0.3);

% Conexoes verticais para o barramento
\draw[line width=1.2pt, draw=buscolor!60] (cpu.south) -- (cpu.south |- 0,-0.3);
\draw[line width=1.2pt, draw=buscolor!60] (gpu.south) -- (gpu.south |- 0,-0.3);
\draw[line width=1.2pt, draw=buscolor!60] (dla.north) -- (dla.north |- 0,-0.3);
\draw[line width=1.2pt, draw=buscolor!60] (video.north) -- (video.north |- 0,-0.3);
\draw[line width=1.2pt, draw=buscolor!60] (vision.north) -- (vision.north |- 0,-0.3);

% ============================================================================
% Conexoes - Blocos para I/O
% ============================================================================
% Linha de barramento I/O
\draw[line width=1.5pt, draw=iocolor!60] (-5, -2.8) -- (5, -2.8);

% Conexoes para I/O
\draw[line width=0.8pt, draw=iocolor!60] (dla.south) -- (dla.south |- 0,-2.8);
\draw[line width=0.8pt, draw=iocolor!60] (video.south) -- (video.south |- 0,-2.8);
\draw[line width=0.8pt, draw=iocolor!60] (vision.south) -- (vision.south |- 0,-2.8);

\draw[line width=0.8pt, draw=iocolor!60] (csi.north) -- (csi.north |- 0,-2.8);
\draw[line width=0.8pt, draw=iocolor!60] (usb.north) -- (usb.north |- 0,-2.8);
\draw[line width=0.8pt, draw=iocolor!60] (pcie.north) -- (pcie.north |- 0,-2.8);
\draw[line width=0.8pt, draw=iocolor!60] (eth.north) -- (eth.north |- 0,-2.8);
\draw[line width=0.8pt, draw=iocolor!60] (gpio.north) -- (gpio.north |- 0,-2.8);

% ============================================================================
% Labels para barramentos
% ============================================================================
\node[labelstyle, anchor=east] at (-5.2, -0.3) {Barramento};
\node[labelstyle, anchor=east] at (-5.2, -0.6) {de Sistema};

\node[labelstyle, anchor=east] at (-5.2, -2.8) {Barramento};
\node[labelstyle, anchor=east] at (-5.2, -3.1) {de I/O};

% ============================================================================
% Titulo do diagrama
% ============================================================================
\node[font=\sffamily\large\bfseries, anchor=south] at (0, 5.3) {
    Arquitetura NVIDIA Jetson AGX Xavier
};

% ============================================================================
% Legenda de potencia
% ============================================================================
\node[font=\sffamily\scriptsize, text=black!60, anchor=north west] at (-5.5, -4.8) {
    Consumo: 10W -- 30W (configuravel)
};

\end{tikzpicture}
}
      \end{center}
    \end{column}
    \begin{column}{0.40\textwidth}
      \vspace{4mm}
      {\small\color{labelcolor}\textbf{Especificações:}}

      \vspace{3mm}
      {\footnotesize
      \begin{itemize}
        \setlength{\itemsep}{3pt}
        \item \textbf{GPU Volta}: 512 CUDA cores + 64 Tensor Cores
        \item \textbf{CPU}: 8x ARM Carmel @ 2,26\,GHz
        \item \textbf{Memória}: 32\,GB LPDDR4x unificada (136,5\,GB/s)
        \item \textbf{IA}: 32 TOPS (INT8)
        \item \textbf{DLA}: 2x aceleradores dedicados
        \item \textbf{Consumo}: 10--30\,W configurável
      \end{itemize}
      }

      \vspace{4mm}
      {\footnotesize\color{stagecolor}\textbf{Escolha:} equilíbrio entre desempenho e eficiência energética para operação contínua.}
    \end{column}
  \end{columns}
\end{frame}

% --- Slide: DeepStream Framework ---
\begin{frame}{NVIDIA DeepStream SDK}
  \vspace{-2mm}
  {\small\color{sublabelcolor}%
    \textit{Framework} para análise de vídeo inteligente com aceleração por \textit{hardware}}

  \vfill
  \begin{center}
    \resizebox{0.95\textwidth}{!}{%
    \begin{tikzpicture}[
      node distance=6mm and 8mm,
      pipeblock/.style={
        rectangle, rounded corners=4pt,
        minimum width=22mm, minimum height=12mm,
        text centered, font=\small\sffamily,
        text=white, draw=none
      },
      arrow/.style={-{Stealth[length=2mm, width=1.5mm]}, thick, arrowcolor},
      grouplabel/.style={font=\footnotesize\sffamily\bfseries, text=labelcolor},
      annot/.style={font=\tiny\sffamily, text=sublabelcolor, align=center},
    ]
      % Pipeline stages
      \node[pipeblock, fill=inputcolor] (src) {Fonte\\de Vídeo};
      \node[pipeblock, fill=backbonecolor, right=of src] (mux) {\texttt{nvstreammux}\\Batch};
      \node[pipeblock, fill=headcolor, right=of mux, minimum width=26mm] (pgie) {PGIE\\\texttt{nvinfer}};
      \node[pipeblock, fill=neckcolor, right=of pgie, minimum width=26mm] (sgie) {SGIEs\\\texttt{nvinfer}};
      \node[pipeblock, fill=cspcolor, right=of sgie] (tracker) {\texttt{nvtracker}\\Rastreamento};
      \node[pipeblock, fill=outputcolor, right=of tracker] (out) {Saída\\Alertas};

      % Arrows
      \draw[arrow] (src) -- (mux);
      \draw[arrow] (mux) -- (pgie);
      \draw[arrow] (pgie) -- (sgie);
      \draw[arrow] (sgie) -- (tracker);
      \draw[arrow] (tracker) -- (out);

      % Annotations below
      \node[annot, below=4mm of src] {RTSP / USB\\Câmeras};
      \node[annot, below=4mm of mux] {Multi-stream\\Multiplexação};
      \node[annot, below=4mm of pgie] {Detecção\\(frame completo)};
      \node[annot, below=4mm of sgie] {Classificação\\(crops)};
      \node[annot, below=4mm of tracker] {DeepSORT\\NvDCF};
      \node[annot, below=4mm of out] {Metadados\\Visualização};

      % Top: acceleration bar
      \draw[thick, skipcolor, decorate, decoration={brace, amplitude=4pt, raise=4pt}]
        (mux.north west) -- (tracker.north east)
        node[midway, above=8pt, font=\footnotesize\sffamily\bfseries, text=skipcolor]
        {Aceleração por GPU / TensorRT};

      % GStreamer base
      \draw[thick, dividercolor, decorate, decoration={brace, amplitude=4pt, raise=4pt, mirror}]
        (src.south west |- 0,-1.5) -- (out.south east |- 0,-1.5)
        node[midway, below=8pt, font=\footnotesize\sffamily\bfseries, text=dividercolor]
        {Pipeline GStreamer};
    \end{tikzpicture}
    }
  \end{center}
  \vfill

  \begin{beamercolorbox}[wd=\textwidth, rounded=true, shadow=false]{block body}
    \centering
    {\footnotesize\color{stagecolor}%
      Inferência em \textbf{cascata nativa} (PGIE $\rightarrow$ SGIEs) --- zero-copy entre CPU/GPU --- suporte a \textbf{múltiplos streams} simultâneos}
  \end{beamercolorbox}
\end{frame}

% --- Slide: Edge Computing vs Cloud ---
\begin{frame}{Edge Computing vs Cloud}
  \vspace{-2mm}
  {\small\color{sublabelcolor}%
    Por que processar na \textbf{borda}? Comparação de paradigmas para análise de vídeo}

  \vfill
  \begin{columns}[c]
    % --- Cloud Column ---
    \begin{column}{0.46\textwidth}
      \begin{center}
        {\normalsize\color{inputcolor}\textbf{Cloud Computing}}
      \end{center}
      \vspace{2mm}

      \begin{tikzpicture}
        \fill[inputcolor!10, rounded corners=5pt] (0,0) rectangle (\textwidth, -4.2);
        \draw[inputcolor, line width=1pt, rounded corners=5pt] (0,0) rectangle (\textwidth, -4.2);

        \node[anchor=north west, font=\footnotesize\sffamily, text=labelcolor, text width=0.9\textwidth]
          at (0.15, -0.25) {
            \begin{itemize}
              \setlength{\itemsep}{3pt}
              \setlength{\leftskip}{-4mm}
              \item[\color{skipcolor}$\boldsymbol\times$] Latência \textbf{alta} (50--200\,ms)
              \item[\color{skipcolor}$\boldsymbol\times$] Requer \textbf{conexão} constante
              \item[\color{headcolor}\checkmark] Computação \textbf{ilimitada}
              \item[\color{skipcolor}$\boldsymbol\times$] Custo com \textbf{largura de banda}
              \item[\color{skipcolor}$\boldsymbol\times$] Riscos de \textbf{privacidade} (LGPD)
            \end{itemize}
          };
      \end{tikzpicture}
    \end{column}

    % --- Separator ---
    \hfill
    {\color{dividercolor}\vrule width 0.5pt}
    \hfill

    % --- Edge Column ---
    \begin{column}{0.46\textwidth}
      \begin{center}
        {\normalsize\color{headcolor}\textbf{Edge Computing}}
      \end{center}
      \vspace{2mm}

      \begin{tikzpicture}
        \fill[headcolor!8, rounded corners=5pt] (0,0) rectangle (\textwidth, -4.2);
        \draw[headcolor, line width=1pt, rounded corners=5pt] (0,0) rectangle (\textwidth, -4.2);

        \node[anchor=north west, font=\footnotesize\sffamily, text=labelcolor, text width=0.9\textwidth]
          at (0.15, -0.25) {
            \begin{itemize}
              \setlength{\itemsep}{3pt}
              \setlength{\leftskip}{-4mm}
              \item[\color{headcolor}\checkmark] Latência \textbf{ultra-baixa} ($<$50\,ms)
              \item[\color{headcolor}\checkmark] Opera \textbf{offline}
              \item[\color{skipcolor}$\boldsymbol\times$] Recursos \textbf{limitados}
              \item[\color{headcolor}\checkmark] Dados processados \textbf{localmente}
              \item[\color{headcolor}\checkmark] Conformidade com \textbf{LGPD}
            \end{itemize}
          };
      \end{tikzpicture}
    \end{column}
  \end{columns}

  \vfill
  \begin{beamercolorbox}[wd=\textwidth, rounded=true, shadow=false]{block body}
    \centering
    {\footnotesize\color{stagecolor}%
      \textbf{Este trabalho}: processamento \textbf{100\% na borda} --- operação autônoma, tempo real, sem dependência de rede}
  \end{beamercolorbox}
\end{frame}
